\chapter{Zeichnen und Modellieren}\label{kap:modellieren}
\sloppy % lange Befehls-Aufforderungen in Schreibmaschinenschrift umbrechen

In diesem Kapitel lernen Sie, wie Sie in \produkt{} Geometrie erzeugen und verändern: Punkte,
Kurven, Flächen und Volumenkörper. Alle Geometrie wird exakt als NURBS (nicht als Dreiecksnetz)
im Geometriekern \emph{Nurbforge} gespeichert; die Darstellung in den Ansichten ist nur eine
Näherung für den Bildschirm. Dadurch bleiben Radien, Kegelschnitte und Freiformflächen auch nach
Transformationen und Booleschen Operationen exakt -- eine wichtige Voraussetzung für die
optischen Berechnungen der Plugins.

Das Kapitel ist in zwei Teile gegliedert. Der erste Teil
(Abschnitte~\ref{sec:befehle-starten} bis~\ref{sec:objektfang}) erklärt die Bedienung, die für
alle Befehle gleich ist: Befehle starten, Punkte und Zahlen eingeben, Befehle abschließen, die
Konstruktionsebene, der Ortho-Modus und der Objektfang. Der zweite Teil
(Abschnitte~\ref{sec:gruppe-draw} bis~\ref{sec:mod-beispiele}) beschreibt jeden Befehl der
Werkzeugspalte, Gruppe für Gruppe, und zeigt an vollständigen Beispielen, wie Sie typische
optische Bauteile -- Linsen, Spiegel, Blenden -- modellieren. Eine knappe alphabetische
Übersicht aller Befehle mit ihren Aufforderungen finden Sie in Anhang~\ref{kap:befehle}, alle
Tastenkürzel in Anhang~\ref{kap:tastenkuerzel}.

\begin{hinweis}
Längen haben in \produkt{} keine eigene Einheit; sie werden in Modelleinheiten eingegeben und
angezeigt. Beim Export nach STEP werden sie als Millimeter geschrieben (siehe
Kapitel~\ref{kap:dateien}). Es empfiehlt sich daher, durchgehend in Millimetern zu modellieren.
Winkel geben Sie immer in Grad ein.
\end{hinweis}

% ============================================================================
\section{Befehle starten}\label{sec:befehle-starten}
\index{Befehl!starten}

Jeder Modellierschritt ist ein \emph{Befehl}: ein kleiner Ablauf, der Sie nacheinander nach
Punkten, Zahlen oder einer Auswahl fragt und am Ende ein Ergebnis erzeugt. Sie können einen
Befehl auf drei Arten starten.

\subsection{Über die Werkzeugspalte}\index{Werkzeugspalte}

Die Werkzeugspalte (Dock \gui{Tools}, standardmäßig links im Hauptfenster) enthält alle Befehle
als Knöpfe, geordnet in Gruppen:

\begin{referenz}{llX}
\toprule
Gruppe & Standard & Befehle \\
\midrule\endhead
\gui{Draw} & aufgeklappt & Point, Line, Tangent line, Perpendicular line, Polyline,
  Rectangle, Polygon \\
\gui{Curves} & aufgeklappt & Circle (centre, radius), Circle (3 points), Arc, Ellipse, Helix,
  Spiral, Curve (control points), Interpolated curve, Sketch (freehand), Curve through points,
  Conic, Parabola, Hyperbola, Offset curve, Fillet curves, Chamfer curves, Blend curves,
  Tween curves \\
\gui{Curves from objects} & zugeklappt & Project curves, Pull curves, Section, Contour,
  Duplicate edge, Duplicate border, Extract isocurve, Extract wireframe, Offset curve on surface,
  Curve Boolean \\
\gui{Surfaces} & zugeklappt & Plane, Extrude, Revolve, Loft, Sweep 1 rail, Sweep 2 rails,
  Coons patch, Gordon surface, Offset surface, Fillet surfaces, Cap planar holes \\
\gui{Solids} & zugeklappt & Box, Sphere, Cylinder, Cone, Torus, Boolean union,
  Boolean difference, Boolean intersection \\
\gui{Transform} & zugeklappt & Move, Copy, Rotate, Scale, Mirror, Linear array, Grid array,
  Polar array \\
\gui{Edit} & zugeklappt & Trim, Split, Join, Explode, Delete \\
\gui{Edit curves} & zugeklappt & Extend curve, Divide, Divide by length, Rebuild curve, Change
  degree, Flip direction, Close curve, Curve seam, Insert knot, Remove knot, Smooth curve, Fair
  curve \\
\gui{Analyze} & zugeklappt & Distance, Angle, Curvature, Area, Volume / mass properties,
  Intersection check \\
\bottomrule
\end{referenz}

Ein Klick auf die Kopfzeile einer Gruppe (Dreieck und Name) klappt sie auf oder zu; mehrere
Gruppen können gleichzeitig offen sein. Welche Gruppen offen sind, merkt sich das Programm für
den nächsten Start. Unter den sieben festen Gruppen erscheint für jedes geladene Plugin eine
eigene Gruppe mit dessen Aktionen (siehe Kapitel~\ref{kap:plugins}).

Die Knöpfe zeigen eine Kurzbeschriftung (etwa \gui{Ln} für Line oder \gui{Ext} für Extrude),
solange für den Befehl kein Symbol hinterlegt ist. Wenn Sie die Maus über einen Knopf halten,
nennt ein Tooltip den vollen Namen und darunter, was Sie tippen müssen, um den Befehl über die
Tastatur zu starten, zum Beispiel \enquote{Line} und \enquote{Type: l + Enter}; die Statusleiste
zeigt den Namen ebenfalls. Die Anzahl der Knopfspalten
richtet sich nach der Breite der Werkzeugspalte -- ziehen Sie den Rand des Docks breiter, um
mehr Knöpfe nebeneinander zu sehen. Die Kurzbeschriftungen aller Knöpfe stehen in der Tabelle
in Anhang~\ref{kap:befehle}.

\subsection{Durch Eintippen des Namens}\index{Befehlszeile}\index{Alias}

Jeder Befehl hat einen Namen und meist einen oder mehrere kurze \emph{Aliasse}. Tippen Sie den
Namen oder Alias in die Befehlszeile am unteren Fensterrand und drücken Sie \keys{Enter}. Groß-
und Kleinschreibung spielt keine Rolle: \code{line}, \code{Line}, \code{LI} und \code{l}
starten alle den Befehl Line.

Sie müssen die Befehlszeile dafür nicht anklicken: Sobald Sie in einer Ansicht tippen, gehen
die Zeichen automatisch in die Befehlszeile, und die Befehlszeile erhält den Tastaturfokus.
Ein typischer Ablauf ist also: Maus in die Ansicht, \code{c} tippen, \keys{Enter} -- und der
Kreisbefehl läuft.

Die wichtigsten Aliasse:

\begin{referenz}{lll@{\hspace{2em}}lll}
\toprule
Alias & Befehl & & Alias & Befehl & \\
\midrule\endhead
\code{pt} & Point & & \code{b} & Box & \\
\code{l} & Line & & \code{sph} & Sphere & \\
\code{pl} & Polyline & & \code{cyl} & Cylinder & \\
\code{rec} & Rectangle & & \code{cone} & Cone & \\
\code{pol} & Polygon & & \code{tor} & Torus & \\
\code{c} & Circle & & \code{bu}, \code{bd}, \code{bi} & Booleans & \\
\code{c3} & Circle (3 points) & & \code{m} & Move & \\
\code{a} & Arc & & \code{co} & Copy & \\
\code{h} & Helix & & \code{ro} & Rotate & \\
\code{sp}, \code{interpcrv} & Interpolated curve & & \code{sc} & Scale & \\
\code{crv} & Curve (control points) & & & & \\
\code{el} & Ellipse & & \code{spi} & Spiral & \\
\code{sk} & Sketch & & \code{crvpts} & Curve through points & \\
\code{conic} & Conic & & \code{para}, \code{hyp} & Parabola, Hyperbola & \\
\code{blend} & Blend curves & & \code{cha} & Chamfer curves & \\
\code{tan} & Tangent line & & \code{perp} & Perpendicular line & \\
\code{project}, \code{pull} & Project / Pull curves & & \code{section}, \code{contour} & Section, Contour & \\
\code{dupedge}, \code{dupborder} & Duplicate edge / border & & \code{iso} & Extract isocurve & \\
\code{offsetcrvonsrf} & Offset curve on surface & & \code{crvboolean} & Curve Boolean & \\
\code{extendcrv} & Extend curve & & \code{divide}, \code{divlen} & Divide (by length) & \\
\code{rebuild} & Rebuild curve & & \code{changedegree} & Change degree & \\
\code{flip} & Flip direction & & \code{closecrv}, \code{crvseam} & Close curve, Curve seam & \\
\code{insertknot}, \code{removeknot} & Knots & & \code{smooth}, \code{fair} & Smooth, Fair & \\
\code{tween} & Tween curves & & \code{filletcorners} & Fillet curves & \\
\code{offset} & Offset curve & & \code{mi} & Mirror & \\
\code{fillet} & Fillet curves & & \code{array} & Linear array & \\
\code{ext} & Extrude & & \code{arraypolar} & Polar array & \\
\code{rev} & Revolve & & \code{j} & Join & \\
\code{offsetsrf} & Offset surface & & \code{x} & Explode & \\
\code{filletsrf} & Fillet surfaces & & \code{tr} & Trim & \\
\code{network} & Gordon surface & & \code{dist} & Distance & \\
\code{volume} & Volume & & \code{ang} & Angle & \\
\code{intersect} & Intersect & & \code{curv} & Curvature & \\
\bottomrule
\end{referenz}

Die vollständige Liste mit allen Namen und Aliassen enthält Anhang~\ref{kap:befehle}.

\begin{warnung}
Verwechseln Sie nicht \code{intersection} und \code{intersect}: \code{intersection} (Alias
\code{bi}) ist die \emph{Boolesche Schnittmenge} von Volumenkörpern, \code{intersect} (Name
\code{intersection\_check}) berechnet die \emph{Schnittpunkte und Schnittkurven} von Kurven und
Flächen (Werkzeug \gui{Intersection check} in der Gruppe Analyze).
\end{warnung}

\subsection{Aliasse sind keine Tastenkürzel}\index{Tastenkürzel!Befehle}

Die Befehle haben keine eigenen Tastenkürzel: Ein Alias wie \code{l} startet einen Befehl erst
mit \keys{Enter}. Der Tooltip jedes Werkzeugs nennt den kürzesten Alias, etwa \enquote{Type: l +
Enter} bei \gui{Line}, \enquote{Type: c + Enter} bei \gui{Circle (centre, radius)} und
\enquote{Type: m + Enter} bei \gui{Move}. Ein echtes Tastenkürzel hat nur das Werkzeug
\gui{Delete}: Es löscht die Auswahl wie \keys{Entf} (siehe Abschnitt~\ref{sec:gruppe-edit}); es
ist kein Befehl der Befehlszeile und hat daher keinen Alias.

\begin{hinweis}
Einige Tasten sind in den Ansichten mit anderen Funktionen belegt und gelangen dann nicht in
die Befehlszeile: \keys{Z} leitet die Tastenfolgen \keys{Z}, \keys{E} (Zoom extents) und
\keys{Z}, \keys{S} (Zoom selected) ein, und die \keys{Leertaste} maximiert die aktive Ansicht.
\keys{H} blendet die Auswahl aus (\gui{Hide}) -- aber nur, wenn etwas ausgewählt ist; ist
nichts ausgewählt, wird \keys{H} wie jeder andere Buchstabe in die Befehlszeile geschrieben.
So starten Sie die Helix auch aus einer Ansicht heraus mit \code{h} und \keys{Enter}, solange
die Auswahl leer ist. In der Befehlszeile selbst werden alle Zeichen ganz normal eingegeben.
\end{hinweis}

\subsection{Den letzten Befehl wiederholen}\index{Befehl!wiederholen}

Drücken Sie \keys{Enter}, während kein Befehl läuft und die Befehlszeile leer ist, oder klicken
Sie ohne laufenden Befehl mit der rechten Maustaste in eine Ansicht (\maus{Rechtsklick ohne
Ziehen}), so wird der zuletzt gestartete Befehl erneut gestartet. Das ist praktisch, wenn Sie mehrere Linien,
Kreise oder Verschiebungen nacheinander ausführen. Wiederholt wird der Befehl mit seinem
Ausgangszustand (etwa wieder mit 6 Ecken beim Polygon); auswahlbasierte Befehle verwenden die
\emph{dann aktuelle} Auswahl.

\subsection{Einen Befehl während eines anderen starten}

Tippen Sie den Namen eines Befehls ein, während ein anderer noch läuft, so wird der laufende
Befehl abgebrochen und der neue gestartet. Bereits erzeugte Zwischenergebnisse (etwa die schon
gesetzten Punkte des Befehls Point) bleiben erhalten. Ist der Name unbekannt, erscheint in der
Statusleiste \code{Unknown command '...'}, und der laufende Befehl bleibt unverändert aktiv.

% ============================================================================
\section{Die Befehlszeile und ihre Meldungen}\label{sec:mod-befehlszeile}

Die Befehlszeile ist das Eingabefeld am unteren Rand des Hauptfensters. Solange kein Befehl
läuft, steht darin grau der Platzhalter \code{Command:}. Läuft ein Befehl, zeigt sie stattdessen
dessen \emph{Aufforderung} (Prompt), die genau sagt, welche Eingabe als Nächstes erwartet wird,
zum Beispiel:

\begin{lstlisting}
Line: start point
Line: end point (or length)
Polygon (6 sides): centre (type a number to change the sides)
\end{lstlisting}

Ist der Ortho-Modus eingeschaltet, wird an die Aufforderung \code{[Ortho]} angehängt. Lesen Sie
die Aufforderung im Zweifel immer zuerst -- sie verrät auch, ob eine Zahl erlaubt ist (\enquote{point
or number}), ob \keys{Enter} eine Bedeutung hat (\enquote{Enter: done}, \enquote{Enter: 360})
und welche Voreinstellung gerade gilt (\enquote{6 sides}, \enquote{5 turns}).

Während Sie die Maus bewegen, zeigt eine orange \emph{Gummiband-Vorschau}\index{Gummiband} in
allen Ansichten, wie das Ergebnis aussehen würde: die Linie bis zum Cursor, den Kreis mit dem
aktuellen Radius, die verschobenen Objekte und so weiter.

Ergebnisse und Fehler meldet das Programm in der Statusleiste links unten:

\begin{referenz}{>{\raggedright\arraybackslash}p{0.38\linewidth}X}
\toprule
Meldung & Bedeutung \\
\midrule\endhead
\code{<Befehl>: <Grund>} & Die Eingabe wurde abgelehnt oder der Geometriekern konnte das
  Ergebnis nicht berechnen; der Text nach dem Doppelpunkt nennt den Grund, zum Beispiel
  \code{Fillet: radius 5 too large for these curves} oder \code{Split: the point is at an end
  of the curve}. Ein Befehl, der auf eine Eingabe wartet, bleibt im selben Schritt; geben Sie
  einen anderen Wert ein. Befehle, die sofort auf der Auswahl arbeiten (etwa Join, Explode,
  Cap), enden ohne Ergebnis. \\
\code{Input not accepted -- <Aufforderung>} & Die Eingabe passt nicht zum aktuellen Schritt,
  und der Befehl nennt keinen eigenen Grund -- zum Beispiel ein Endpunkt, der mit dem
  Startpunkt zusammenfällt, ein Radius 0 oder drei Punkte auf einer Geraden beim Kreis durch
  drei Punkte. Der Befehl bleibt im selben Schritt; geben Sie einen anderen Wert ein. \\
\code{Unknown command '...'} & Der eingetippte Name ist weder Befehl noch Alias. \\
\code{No command is waiting for a point} & Sie haben Koordinaten eingegeben, ohne dass ein
  Befehl läuft. \\
\code{A number is not expected here} & Eine Zahl wurde eingegeben, der aktuelle Schritt kann
  aber weder eine Zahl noch eine Länge verwenden. \\
\code{relative input needs a previous point} & Eine Eingabe mit \code{@} ist erst möglich,
  wenn es einen Bezugspunkt gibt (Abschnitt~\ref{sec:punkteingabe}). \\
\code{polar input: distance<angle} & Die Polar-Eingabe ist nicht in der Form
  \code{Abstand<Winkel}. \\
\code{expected x,y or x,y,z} & Eine Koordinateneingabe hat nicht zwei oder drei Werte. \\
\code{'...' is not a number} & Ein Teil einer Koordinateneingabe ist keine Zahl. \\
\code{<Befehl> cancelled} & Der Befehl wurde abgebrochen, ohne etwas zu erzeugen. \\
\code{<Befehl>: select objects first} & Ein auswahlbasierter Befehl wurde ohne Auswahl
  gestartet (Abschnitt~\ref{sec:auswahlbefehle}). \\
\code{Boolean ... failed: ...} & Eine Boolesche Operation ist gescheitert; der Text nach dem
  Doppelpunkt nennt den Grund (Abschnitt~\ref{sec:boolesche}). \\
\bottomrule
\end{referenz}

Messergebnisse (Distance, Angle, Area, Volume, Curvature) erscheinen ebenfalls in der
Statusleiste. Das Ergebnis von Distance bleibt dort stehen, bis eine neue Meldung kommt; die
übrigen Meldungen verschwinden nach einigen Sekunden.

% ============================================================================
\section{Punkte eingeben}\label{sec:punkteingabe}
\index{Punkteingabe}\index{Koordinaten}

Fast jeder Befehl fragt nach Punkten. Sie können jeden Punkt wahlweise mit der Maus setzen oder
als Koordinaten eintippen; beides lässt sich innerhalb eines Befehls beliebig mischen.

\subsection{Mit der Maus}

Während ein Befehl einen Punkt erwartet, wird der Cursor in den Ansichten zum Fadenkreuz.
\maus{Klicken} Sie mit der linken Maustaste, um den Punkt zu setzen. Welcher Punkt genau
entsteht, hängt davon ab, ob der Objektfang greift:

\begin{itemize}
\item \textbf{Mit Fang:} Liegt ein Fangpunkt in der Nähe des Cursors -- etwa ein Kurvenende oder
  ein Kreismittelpunkt --, zeigt ein grünes Symbol mit Beschriftung diesen Punkt an, und der
  Klick übernimmt seine exakten 3D-Koordinaten (Abschnitt~\ref{sec:objektfang}).
\item \textbf{Ohne Fang:} Sonst entsteht der Punkt dort, wo der Sehstrahl unter dem Cursor die
  \emph{Konstruktionsebene} der Ansicht trifft (Abschnitt~\ref{sec:mod-konstruktionsebene}). In der
  Ansicht Top ist das die XY-Ebene bei $z=0$, in Front die XZ-Ebene bei $y=0$ usw.
\end{itemize}

Die Statusleiste zeigt laufend die Koordinaten des Punktes, der bei einem Klick entstehen würde
(\code{X ...   Y ...   Z ...} mit drei Nachkommastellen) -- beim Fangen die Koordinaten des
Fangpunkts.

Während der Punkteingabe bleibt die Navigation in den Ansichten möglich: Mausrad zum Zoomen,
mittlere Maustaste zum Verschieben, \maus{rechts ziehen} zum Drehen bzw.\ Verschieben (siehe
Kapitel~\ref{kap:navigation}). Ein \maus{Rechtsklick ohne Ziehen} beendet dagegen den Befehl wie
\keys{Enter} (Abschnitt~\ref{sec:abschliessen}). Ein Ziehen mit gedrückter linker Taste zieht
während eines Befehls \emph{kein} Auswahlrechteck auf; der Punkt entsteht dort, wo Sie die Taste
loslassen.

\subsection{Über die Tastatur}\label{sec:koordinateneingabe}

Tippen Sie Koordinaten ein und bestätigen Sie mit \keys{Enter}. Dezimaltrennzeichen ist immer
der \emph{Punkt}, die Koordinaten werden durch \emph{Kommas} getrennt; Leerzeichen sind erlaubt
und werden ignoriert. Wissenschaftliche Schreibweise wie \code{1.5e-3} ist zulässig.

\begin{referenz}{>{\raggedright\arraybackslash}p{0.22\linewidth}X}
\toprule
Eingabe & Bedeutung \\
\midrule\endhead
\code{x,y,z} & absolute Weltkoordinaten, z.\,B.\ \code{10,0,-5} \\
\code{x,y} & Koordinaten in der Konstruktionsebene der Ansicht (siehe unten), gemessen vom
  Weltursprung \\
\code{@dx,dy,dz} & relativ zum Bezugspunkt, Versatz in Weltkoordinaten \\
\code{@dx,dy} & relativ zum Bezugspunkt, Versatz in der Konstruktionsebene \\
\code{d<a} & polar in der Konstruktionsebene: Abstand $d$ vom Ursprung unter dem Winkel $a$ (in
  Grad, gemessen von der ersten Ebenenachse $u$ zur zweiten Achse $v$) \\
\code{@d<a} & polar relativ zum Bezugspunkt \\
\bottomrule
\end{referenz}

Was \enquote{x} und \enquote{y} bei der Eingabe mit zwei Werten bedeuten, hängt von der Ansicht
ab, in der Sie arbeiten: Beim Start eines Befehls ist das die \emph{aktive} Ansicht (die Sie
zuletzt angeklickt haben), danach die Ansicht, über der die Maus während des Befehls zuletzt
war:

\begin{referenz}{lllX}
\toprule
Ansicht & 1.\ Wert ($u$) & 2.\ Wert ($v$) & Beispiel \code{3,4} ergibt \\
\midrule\endhead
Top (XY) & $x$ & $y$ & $(3,\,4,\,0)$ \\
Perspective & $x$ & $y$ & $(3,\,4,\,0)$ \\
Front (XZ) & $x$ & $z$ & $(3,\,0,\,4)$ \\
Left (YZ) & $-y$ & $z$ & $(0,\,-3,\,4)$ \\
\bottomrule
\end{referenz}

\begin{tipp}
Bewegen Sie die Maus nach dem Start des Befehls kurz über die gewünschte Ansicht, \emph{bevor}
Sie Koordinaten oder Zahlen eintippen. Die Konstruktionsebene, auf die sich \code{x,y},
\code{@dx,dy}, \code{d<a}, eine eingetippte Extrude-Distanz und viele Befehle (Rechteck, Kreis,
Zylinder, Rotation \ldots) beziehen, ist die der Ansicht, über der die Maus während des Befehls
zuletzt war -- oder, solange die Maus noch über keiner Ansicht war, die der aktiven Ansicht. Mit
drei Werten (\code{x,y,z}) sind Sie dagegen unabhängig von der Ansicht.
\end{tipp}

\paragraph{Der Bezugspunkt.}\index{Bezugspunkt}
Relative Eingaben (\code{@\ldots}), eingetippte Längen und der Ortho-Modus beziehen sich auf den
\emph{Bezugspunkt} des laufenden Befehls. In der Regel ist das der zuletzt eingegebene Punkt
(bei Line der Startpunkt, bei Polyline der letzte Eckpunkt). Einige Befehle verwenden einen
festen Bezugspunkt: Circle, Polygon, Sphere, Cylinder, Cone und Torus den Mittelpunkt, Rotate
und Polar array das Zentrum, Scale den Basispunkt, Angle den Scheitel und Grid array den
Basispunkt. Bevor der erste Punkt eines Befehls gesetzt ist, gibt es keinen Bezugspunkt; eine
Eingabe mit \code{@} wird dann mit der Meldung \code{relative input needs a previous point}
abgelehnt.

\paragraph{Beispiele.}
Ein Quadrat mit 20 Einheiten Kantenlänge als Polylinie in der Ansicht Top:
\begin{lstlisting}
pl            (Enter)  Polyline: start point
0,0           (Enter)  Polyline: next point (Enter: done, snap to the first point: close)
@20,0         (Enter)
@0,20         (Enter)
@-20,0        (Enter)
0,0           (Enter)  -> der Startpunkt schliesst die Polylinie
\end{lstlisting}
Eine Linie der Länge 50 unter \SI{30}{\degree} vom Punkt $(10,10)$ aus:
\begin{lstlisting}
l             (Enter)  Line: start point
10,10         (Enter)  Line: end point (or length)
@50<30        (Enter)
\end{lstlisting}

\subsection{Zahlen eingeben}\label{sec:zahleneingabe}\index{Zahleneingabe}

Viele Schritte erlauben statt eines Punkts eine Zahl -- die Aufforderung sagt das mit
\enquote{(point or number)}, \enquote{or radius}, \enquote{or length} oder \enquote{type a
number}. Eine eingetippte Zahl wird so verwendet:

\begin{enumerate}
\item Erwartet der aktuelle Schritt ausdrücklich eine Zahl (Radius, Höhe, Anzahl der Ecken,
  Windungen, Winkel, Faktor, Abstand \ldots), wird sie genau dafür verwendet.
\item Andernfalls wird die Zahl als \emph{Länge} interpretiert: Der Punkt entsteht im Abstand
  dieser Länge vom Bezugspunkt, in Richtung des Mauszeigers. So zeichnen Sie eine Linie der
  Länge 25, indem Sie den Startpunkt setzen, die Maus in die gewünschte Richtung bewegen und
  \code{25} eintippen. Zusammen mit dem Ortho-Modus (\keys{F8}) ergibt das schnell achsparallele
  Strecken genau bekannter Länge. Die Länge muss positiv sein.
\item Ist beides nicht möglich (kein Bezugspunkt, Maus außerhalb der Ansichten), erscheint
  \code{A number is not expected here}.
\end{enumerate}

Einige Befehle verwenden Zahlen am Anfang als \emph{Option}: Bevor Sie den ersten Punkt setzen,
ändert eine Zahl beim Polygon die Eckenzahl, bei der Helix die Windungszahl und bei den Reihen
(Arrays) die Anzahl der Kopien. Die Aufforderung zeigt den jeweils gültigen Wert in Klammern an
und kündigt die Option mit \enquote{type a number to change \ldots} an. Weitere Optionen in
Form von Schlüsselwörtern gibt es nicht.

% ============================================================================
\section{Befehle abschließen und abbrechen}\label{sec:abschliessen}
\index{Befehl!abschließen}\index{Befehl!abbrechen}\index{Rückgängig}

\begin{referenz}{>{\raggedright\arraybackslash}p{0.38\linewidth}X}
\toprule
Eingabe & Wirkung \\
\midrule\endhead
\keys{Enter} (leere Befehlszeile) oder \maus{Rechtsklick} (ohne Ziehen) in einer Ansicht &
  beendet den laufenden Befehl. Befehle mit offener Punktanzahl (Point, Polyline, Spline, Copy)
  erzeugen dabei ihr Ergebnis; Revolve und Polar array verwenden \SI{360}{\degree}. Läuft kein
  Befehl, wiederholen \keys{Enter} und der Rechtsklick den letzten Befehl. Ein Ziehen mit der
  rechten Maustaste dreht bzw.\ verschiebt nur die Ansicht. \\
\keys{Esc} & bricht den laufenden Befehl ab (\code{<Befehl> cancelled}); was der Befehl noch
  nicht erzeugt hat, wird verworfen. Läuft kein Befehl, bricht \keys{Esc} einen laufenden
  Plugin-Auftrag ab, und erst wenn auch der nicht läuft, hebt es die Auswahl auf (Menü
  \menu{Edit,Select None / Cancel}). \\
\keys{Strg+Z} & macht das letzte Ergebnis rückgängig (\menu{Edit,Undo}). \\
\bottomrule
\end{referenz}

Jedes Ergebnis eines Befehls ist \emph{ein} Rückgängig-Schritt, der im Menü \menu{Edit} mit dem
Namen des Befehls erscheint (etwa \enquote{Undo Line}). Das gilt auch für Befehle, die
Eingangsobjekte entfernen oder ersetzen, etwa Boolesche Operationen oder Join: Ein
\keys{Strg+Z} stellt den Zustand vor dem Befehl vollständig wieder her. Bei Point und Copy, die
mehrere Ergebnisse nacheinander erzeugen, ist jeder gesetzte Punkt bzw.\ jede Kopie ein eigener
Schritt.

Erzeugt ein Befehl beim Beenden mit \keys{Enter} nichts (etwa eine Polylinie mit nur einem
Punkt), endet er wie bei \keys{Esc} mit der Meldung \code{<Befehl> cancelled}. Bei Befehlen mit
fester Anzahl von Eingaben (etwa Line, Circle, Box) bricht \keys{Enter} vor der letzten Eingabe
den Befehl also ab.

% ============================================================================
\section{Befehle auf der Auswahl}\label{sec:auswahlbefehle}
\index{Auswahl!Reihenfolge}

Viele Befehle arbeiten nicht mit neu eingegebenen Punkten, sondern mit vorhandenen Objekten:
Transformationen, Boolesche Operationen, Flächen aus Kurven, Join, Trim, die Messungen und
weitere. Für diese Befehle gilt:

\begin{itemize}
\item \textbf{Erst auswählen, dann den Befehl starten.} Der Befehl übernimmt beim Start die
  aktuelle Auswahl. Während der Befehl läuft, setzt ein Klick in die Ansicht Punkte und wählt
  keine Objekte. Ist nichts (Passendes) ausgewählt, meldet die Statusleiste, was fehlt -- etwa
  \code{Move: select objects first} oder \code{Loft: select two or more section curves (in
  order)} --, und der Befehl startet nicht.
\item \textbf{Ganze Objekte.} Es zählen nur ganz ausgewählte Objekte. Einzeln ausgewählte
  Flächen, Kanten oder Ecken eines Körpers (Auswahl mit \keys{Strg+Shift}+Klick) werden von den
  Befehlen ignoriert.
\item \textbf{Die Reihenfolge zählt.} Die Objekte werden in der Reihenfolge übergeben, in der
  Sie sie ausgewählt haben. Bei Boolean difference ist das zuerst gewählte Objekt der Körper,
  von dem abgezogen wird; bei Trim und Split das zu schneidende Objekt; bei Loft die Reihenfolge
  der Querschnitte; bei Sweep das Profil vor der Leitkurve. Wählen Sie solche Objekte
  nacheinander mit \maus{Klick} und \keys{Shift}+\maus{Klick} aus, nicht mit einem
  Auswahlrechteck.
\item \textbf{Passende Typen.} Objekte, die für den Befehl nicht passen (etwa Punkte bei Loft),
  werden stillschweigend übergangen.
\end{itemize}

Wie Sie Objekte auswählen, beschreibt Kapitel~\ref{kap:navigation}.

% ============================================================================
\section{Die Konstruktionsebene}\label{sec:mod-konstruktionsebene}
\index{Konstruktionsebene}

Jede der vier Ansichten hat eine \emph{Konstruktionsebene} mit zwei Achsen $u$ und $v$ und der
Normalen $n = u \times v$. Das Raster der Ansicht liegt in dieser Ebene. Sie bestimmt,

\begin{itemize}
\item wo ein Mausklick ohne Fang landet (auf der Ebene),
\item wie zweiwertige Koordinaten \code{x,y}, \code{@dx,dy} und Polarkoordinaten gemeint sind,
\item in welcher Ebene Rechteck, Polygon, Kreis, Plane und die Grundflächen von Box, Zylinder,
  Kegel und Torus liegen,
\item die Richtung der Höhe bei Box, Cylinder und Cone (positive Werte entlang $n$) und einer
  eingetippten Distanz bei Extrude,
\item die Drehachse bei Rotate und Polar array und die Spiegelebene bei Mirror,
\item die Achsen des Ortho-Modus.
\end{itemize}

\begin{referenz}{llllX}
\toprule
Ansicht & Ebene & $u$ & $v$ & Normale $n$ \\
\midrule\endhead
Top (XY) & $z=0$ & $+x$ & $+y$ & $+z$ (zum Betrachter) \\
Perspective & $z=0$ & $+x$ & $+y$ & $+z$ \\
Front (XZ) & $y=0$ & $+x$ & $+z$ & $-y$ (zum Betrachter) \\
Left (YZ) & $x=0$ & $-y$ & $+z$ & $-x$ (zum Betrachter) \\
\bottomrule
\end{referenz}

Alle Konstruktionsebenen gehen durch den Weltursprung. Stellen Sie eine Ansicht über ihr
Ansichtsmenü (Klick auf den Titel der Ansicht) auf eine andere Standardansicht um, übernimmt sie
auch deren Konstruktionsebene. Eigene, frei verschobene oder gedrehte Konstruktionsebenen gibt
es in dieser Version nicht.

\begin{tipp}
Um in einer anderen Höhe zu zeichnen, geben Sie den ersten Punkt mit drei Koordinaten ein: Ein
Kreis mit dem Mittelpunkt \code{0,0,20}, gezeichnet über der Ansicht Top, liegt in der Ebene
$z=20$ -- der Kreis wird in der Ebene durch seinen Mittelpunkt parallel zur Konstruktionsebene
erzeugt. Dasselbe gilt für Polygon, Rechteck (erster Eckpunkt) und die Grundfläche der
Volumenkörper. Alternativ fangen Sie den ersten Punkt an vorhandener Geometrie.
\end{tipp}

% ============================================================================
\section{Ortho-Modus}\label{sec:ortho}
\index{Ortho}

Der Ortho-Modus hilft, Punkte genau achsparallel zum Bezugspunkt zu setzen. Sie schalten ihn mit
\keys{F8} oder mit dem Knopf \gui{Ortho} rechts in der Statusleiste ein und aus; im
eingeschalteten Zustand ist der Knopf gedrückt dargestellt, und die Aufforderung endet mit
\code{[Ortho]}.

Bei eingeschaltetem Ortho wird ein mit der Maus gesetzter Punkt auf diejenige Achse der
Konstruktionsebene durch den Bezugspunkt verschoben ($u$- oder $v$-Richtung), die der
Cursorrichtung am nächsten liegt. In der Ansicht Top entstehen so nur waagrechte und
senkrechte Linien. Das Gummiband zeigt die korrigierte Lage bereits an.

Ortho wirkt nicht
\begin{itemize}
\item auf Punkte, die über den Objektfang gefangen wurden -- ein gefangener Punkt hat immer
  Vorrang und wird nicht verändert;
\item auf eingetippte Koordinaten;
\item solange der Befehl noch keinen Bezugspunkt hat (also beim ersten Punkt).
\end{itemize}
Eine eingetippte Länge (Abschnitt~\ref{sec:zahleneingabe}) nimmt dagegen die Ortho-korrigierte
Richtung zum Cursor.

\begin{tipp}
Wenn ein Punkt mit Ortho nicht dort landet, wo Sie ihn erwarten, hat vermutlich der Objektfang
gegriffen. Halten Sie \keys{Alt} gedrückt, um den Fang für diesen Klick zu unterdrücken, oder
schalten Sie ihn mit \keys{F3} ab.
\end{tipp}

% ============================================================================
\section{Objektfang}\label{sec:objektfang}
\index{Objektfang}\index{Fang|see{Objektfang}}

Der Objektfang (\enquote{magnetischer Cursor}) zieht den Cursor auf charakteristische Punkte
der vorhandenen Geometrie: Enden, Mittelpunkte, Zentren, Schnittpunkte und andere. So setzen Sie
Punkte exakt -- in doppelter Genauigkeit, direkt aus der NURBS-Geometrie berechnet, nicht aus der
Bildschirmdarstellung. Der Fang ist aktiv, solange ein Befehl einen Punkt erwartet, und
funktioniert in allen vier Ansichten und in jedem Darstellungsmodus.

\subsection{Fangsymbole}

Ein grünes Symbol markiert den Fangpunkt; daneben steht in einem dunklen Kästchen die Fangart.
Abbildung~\ref{fig:fangsymbole} zeigt alle Symbole.

\abbildung[0.95\linewidth]{snap-glyphs.png}{Die Fangsymbole mit ihren Beschriftungen, in der
Reihenfolge der Statusleiste: Point, End, Corner, Int, Cen, Quad, Mid, Edge, Near,
Grid.}{fig:fangsymbole}

\subsection{Die Fangarten}

\begin{referenz}{llp{0.17\linewidth}X}
\toprule
Fangart & Kürzel & Symbol & Was gefangen wird \\
\midrule\endhead
Punkt & \gui{Point} & kleines gefülltes Quadrat & Punktobjekte (erzeugt mit Point oder
  Intersect). \\
Ende & \gui{End} & Quadrat & Anfangs- und Endpunkt offener Kurven: Linien, Polylinien,
  Bögen, Splines, Helices. Dazu \emph{jede Ecke} einer Polylinie, eines Rechtecks, eines
  Polygons und einer mit Join verbundenen Kurve -- bei geschlossenen Kurven auch die Ecke am
  Anfang/Ende (Naht). Bei Kurven höheren Grades (Splines, Bögen, verbundene Kurven) zählen nur
  echte Knicke, an denen die Tangente springt. Glatte geschlossene Kurven wie Kreise haben keine
  Enden. Beim Zeichnen einer Polylinie sind auch ihre bereits gesetzten Punkte fangbar, so dass
  Sie sie am Startpunkt schließen können (Abschnitt~\ref{sec:cmd-polyline}). \\
Flächenecke & \gui{Corner} & Quadrat mit Punkt & Ecken von Flächen und Körpern: die
  Endpunkte der Randkanten (B-Rep-Ecken), etwa die Ecken einer Box. \\
Schnittpunkt & \gui{Int} & X & Schnittpunkte zweier Kurven oder Kanten in 3D. In den
  Parallelprojektionen (Front, Left, Top) zusätzlich \emph{scheinbare} Schnittpunkte: Kurven,
  die sich nur in der Ansicht kreuzen, im Raum aber aneinander vorbeilaufen. Der Fang nimmt dann
  den Punkt auf der vorderen Kurve; die Beschriftung lautet \code{Int (apparent)}. \\
Zentrum & \gui{Cen} & Kreis & Mittelpunkte von Kreisen und Bögen (Krümmungsmittelpunkt);
  bei anderen geschlossenen ebenen Kurven, etwa einem Rechteck oder Polygon, der
  Flächenschwerpunkt. \\
Quadrant & \gui{Quad} & Raute & Punkte, in denen Kreise und Bögen die Ebenen durch ihren
  Mittelpunkt parallel zu den Koordinatenebenen schneiden -- bei einem Kreis in der XY-Ebene die
  Punkte bei 0, 90, 180 und \SI{270}{\degree}. Bei Bögen nur die Punkte, die auf dem Bogen
  liegen. \\
Mitte & \gui{Mid} & Dreieck & Punkt auf halber \emph{Bogenlänge} offener Kurven (bei
  Splines also nicht die Parametermitte). Hat eine Kurve Ecken (Polylinie, Rechteck, Polygon,
  verbundene Kurve, auch geschlossen), hat stattdessen \emph{jedes Segment} zwischen zwei Ecken
  seine eigene Mitte. \\
Flächenkante & \gui{Edge} & (Schalter) & Kein eigener Fangpunkt, sondern ein Schalter: Ist er
  an, nehmen auch die Kanten von Flächen und Körpern an den Fangarten Mid, Near, Cen und Quad
  teil -- so fangen Sie zum Beispiel den Mittelpunkt der kreisförmigen Kante einer Bohrung. \\
Nächster & \gui{Near} & Sanduhr & Der Punkt auf einer Kurve oder Kante, der dem Cursor am
  nächsten liegt. Damit setzen Sie Punkte \enquote{irgendwo auf} einer Kurve. \\
Raster & \gui{Grid} & kleines Kreuz & Rasterpunkte der Konstruktionsebene (Schrittweite der
  feinen Rasterlinien). Standardmäßig ausgeschaltet. \\
\bottomrule
\end{referenz}

\enquote{Kurve} heißt hier jede Kurve: Linien, Polylinien, Bögen, Kreise und
Freiform-NURBS-Kurven. Kreise und Bögen werden dabei an ihrer Form erkannt, nicht daran, mit
welchem Befehl sie entstanden sind -- auch ein exakt kreisförmiger Spline hat ein Zentrum.

\subsection{Welcher Fangpunkt gewinnt?}\index{Objektfang!Priorität}

\begin{enumerate}
\item \textbf{Fangradius:} Ein Kandidat zählt nur, wenn er auf dem Bildschirm höchstens
  \num{10} Pixel vom Cursor entfernt ist. Der Fang verhält sich dadurch bei jeder Zoomstufe
  gleich: Zoomen Sie hinein, wenn viele Kandidaten dicht beieinander liegen.
\item \textbf{Priorität:} Liegen mehrere Kandidaten im Fangradius, gewinnt die Fangart mit der
  höheren Priorität:
  \begin{center}
  Point $>$ End $>$ Corner $>$ Int $>$ Cen $>$ Quad $>$ Mid $>$ Edge $>$ Near $>$ Grid
  \end{center}
  Innerhalb derselben Fangart gewinnt der Kandidat, der auf dem Bildschirm am nächsten am
  Cursor liegt, bei Gleichstand der, der näher an der Kamera liegt.
\item \textbf{Verdeckte Geometrie:} In den Darstellungsmodi \gui{Surfaces} und \gui{Shaded}
  werden Kandidaten ignoriert, die hinter einer Fläche liegen. Im Modus \gui{Wireframe} sind
  alle Kandidaten fangbar -- schalten Sie auf Wireframe, um einen Punkt auf der Rückseite eines
  Körpers zu fangen.
\end{enumerate}

\subsection{Zwischen Kandidaten wechseln: Tab}\index{Objektfang!Tab}

Liegen mehrere Kandidaten unter dem Cursor, zeigt die Beschriftung zusätzlich ihre Anzahl, etwa
\code{End  1/3}. Drücken Sie \keys{Tab}, um zum nächsten Kandidaten zu wechseln, und
\keys{Shift+Tab}, um zum vorherigen zurückzugehen. Ein anschließender Klick (ohne die Maus zu
bewegen) setzt den gewählten Kandidaten. Bewegen Sie die Maus, beginnt die Auswahl wieder beim
besten Kandidaten.

\subsection{Fang ein- und ausschalten}\index{Objektfang!ausschalten}

\begin{referenz}{>{\raggedright\arraybackslash}p{0.38\linewidth}X}
\toprule
Bedienung & Wirkung \\
\midrule\endhead
\keys{F3} oder Knopf \gui{Snap} in der Statusleiste & schaltet den gesamten Objektfang ein
  und aus. \\
Knöpfe \gui{Point} \ldots{} \gui{Grid} in der Statusleiste & schalten einzelne Fangarten ein
  und aus. Eingeschaltete Fangarten sind gedrückt dargestellt. \\
\keys{Alt} gedrückt halten & unterdrückt den Fang, solange die Taste gehalten wird. Ein
  \keys{Alt}+\maus{Klick} ohne Mausbewegung setzt einen freien Punkt auf der
  Konstruktionsebene; \keys{Alt}+\maus{Ziehen} dreht bzw.\ verschiebt dagegen weiterhin die
  Ansicht. \\
\bottomrule
\end{referenz}

Standardmäßig sind alle Fangarten außer \gui{Grid} eingeschaltet. Der Zustand des Hauptschalters
und der einzelnen Fangarten bleibt über einen Neustart des Programms erhalten.

\begin{tipp}
\begin{itemize}
\item Schalten Sie \gui{Grid} ein, wenn Sie grob auf glatten Werten
  skizzieren wollen; das Raster passt seine Schrittweite der Zoomstufe an (Zehnerpotenzen).
\item Wenn \gui{Near} ständig andere Fangarten überdeckt, ist meist nichts Besseres im
  Fangradius; zoomen Sie hinein oder schalten Sie \gui{Near} vorübergehend ab.
\item Um ein Zentrum zu fangen, bewegen Sie den Cursor auf den Mittelpunkt des Kreises, nicht
  auf die Kreislinie -- der Fangpunkt muss im Fangradius liegen.
\item Eine Polylinie, ein Rechteck und ein Polygon sind jeweils \emph{eine} Kurve, dennoch ist
  jede ihrer Ecken ein \gui{End} und jedes Segment hat ein \gui{Mid}. Sie müssen die Kurve
  dafür nicht mit \gui{Explode} zerlegen.
\end{itemize}
\end{tipp}

% ============================================================================
\section{Zeichnen (Gruppe Draw)}\label{sec:gruppe-draw}

Die Befehle dieser Gruppe erzeugen Punkte und gerade Kurven. Alle Kurven sind NURBS-Kurven;
Polylinie, Rechteck und Polygon sind jeweils \emph{eine} Kurve vom Grad~1 mit Knicken an den
Ecken. Jede Ecke ist ein Fangpunkt \gui{End}, jedes Segment hat ein \gui{Mid}; mit Explode
zerlegen Sie die Kurve bei Bedarf in einzelne Strecken.

\subsection{Point -- Punkt}\label{sec:cmd-point}
\cmd{point}, Aliasse \code{pt}, \code{po}. Setzt Punktobjekte.
\begin{schritte}
\item Starten Sie \gui{Point}. Aufforderung: \code{Point: location (Enter / right-click: done)}.
\item Klicken oder tippen Sie beliebig viele Punkte; jeder Punkt wird sofort erzeugt.
\item Beenden Sie mit \keys{Enter} oder \maus{Rechtsklick}.
\end{schritte}
Punktobjekte dienen als Hilfspunkte für den Fang (Fangart \gui{Point}) und als Markierungen.

\subsection{Line -- Linie}\label{sec:cmd-line}
\cmd{line}, Aliasse \code{l}, \code{li}. Eine Strecke zwischen zwei Punkten.
\begin{schritte}
\item \code{Line: start point} -- Startpunkt setzen.
\item \code{Line: end point (or length)} -- Endpunkt setzen, oder eine Länge eintippen: Die
  Linie verläuft dann vom Startpunkt in Richtung Cursor.
\end{schritte}
Der Befehl endet nach einer Linie; mit \keys{Enter} starten Sie die nächste. Fallen Start- und
Endpunkt zusammen, wird der Endpunkt abgelehnt.

\subsection{Tangent line -- Tangente}\label{sec:cmd-tangent}
\cmd{tangent}, Alias \code{tan}. Eine Strecke, die eine oder zwei ausgewählte ebene Kurven
berührt, in der Konstruktionsebene der Ansicht.
\begin{itemize}
\item \textbf{Eine Kurve}: \code{Tangent: start point}, dann \code{Tangent: click near the
  touching point you want} -- vom Startpunkt gibt es meist mehrere Tangenten (an einen Kreis
  zwei); die Vorschau zeigt jeweils die, deren Berührpunkt dem Cursor am nächsten liegt.
\item \textbf{Zwei Kurven}: \code{Tangent: click near the tangent you want} -- eine gemeinsame
  Tangente beider Kurven (an zwei Kreise bis zu vier), die nächste zum Cursor.
\end{itemize}

\subsection{Perpendicular line -- Lot auf eine Kurve}\label{sec:cmd-perpendicular}
\cmd{perpendicular}, Alias \code{perp}. Wählen Sie eine Kurve aus und klicken Sie Punkte: Jeder
Klick erzeugt die Strecke von dem Punkt zu seinem Fußpunkt auf der Kurve (dem nächsten Punkt;
die Strecke trifft die Kurve im rechten Winkel). \keys{Enter} beendet den Befehl.

\subsection{Polyline -- Polylinie}\label{sec:cmd-polyline}
\cmd{polyline}, Alias \code{pl}. Ein Linienzug aus geraden Segmenten.
\begin{schritte}
\item \code{Polyline: start point}.
\item \code{Polyline: next point (Enter: done, snap to the first point: close)} -- beliebig
  viele weitere Punkte.
\item Abschluss: \keys{Enter} oder \maus{Rechtsklick} erzeugt eine \emph{offene} Polylinie (ab
  zwei Punkten). Ein Klick auf den \emph{Startpunkt} (ab dem vierten Punkt, also mit mindestens
  drei vorhandenen) schließt sie und beendet den Befehl. Der Punkt muss exakt mit dem
  Startpunkt übereinstimmen: Die bereits gesetzten Punkte der entstehenden Polylinie sind
  während des Zeichnens als \gui{End} fangbar -- bewegen Sie den Cursor auf den Startpunkt, bis
  das Fangsymbol erscheint, und klicken Sie. Alternativ tippen Sie seine Koordinaten erneut
  ein.
\end{schritte}
Zwei aufeinanderfolgende gleiche Punkte werden abgelehnt. Eine geschlossene Polylinie hat einen
Zentrumsfangpunkt (Flächenschwerpunkt) und lässt sich zu einem Prisma extrudieren und deckeln
(Beispiel in Abschnitt~\ref{sec:bsp-prisma}).

\subsection{Rectangle -- Rechteck}\label{sec:cmd-rectangle}
\cmd{rectangle}, Aliasse \code{rec}, \code{rect}. Ein geschlossenes Rechteck, dessen Kanten
parallel zu den Achsen der Konstruktionsebene liegen.
\begin{schritte}
\item \code{Rectangle: first corner}.
\item \code{Rectangle: opposite corner} -- die diagonal gegenüberliegende Ecke.
\end{schritte}
Das Rechteck liegt in der Ebene durch die erste Ecke parallel zur Konstruktionsebene der Ansicht,
in der Sie die erste Ecke gesetzt haben. Für ein Rechteck mit bekannten Seitenlängen tippen Sie
die zweite Ecke relativ ein, etwa \code{@40,25}. Hat das Rechteck die Breite oder Höhe 0, wird
die Ecke abgelehnt.

\subsection{Polygon -- regelmäßiges Vieleck}\label{sec:cmd-polygon}
\cmd{polygon}, Alias \code{pol}. Ein regelmäßiges, geschlossenes Vieleck in der
Konstruktionsebene; standardmäßig ein Sechseck.
\begin{schritte}
\item \code{Polygon (6 sides): centre (type a number to change the sides)} -- tippen Sie
  optional zuerst die Eckenzahl (ganze Zahl von 3 bis 1000), dann setzen Sie den Mittelpunkt.
\item \code{Polygon (6 sides): corner point or radius} -- klicken Sie eine Ecke (bestimmt Größe
  und Drehlage) oder tippen Sie den Umkreisradius. Bei einem getippten Radius liegt die erste
  Ecke in $u$-Richtung der Konstruktionsebene.
\end{schritte}

% ============================================================================
\section{Kurven (Gruppe Curves)}\label{sec:gruppe-curves}

\subsection{Circle (centre, radius) -- Kreis um Mittelpunkt}\label{sec:cmd-circle}
\cmd{circle}, Aliasse \code{c}, \code{ci}.
\begin{schritte}
\item \code{Circle: centre}.
\item \code{Circle: radius (point or number)} -- klicken Sie einen Punkt (der Radius ist sein
  Abstand vom Mittelpunkt, gemessen in der Konstruktionsebene) oder tippen Sie den Radius.
\end{schritte}
Der Kreis liegt in der Ebene durch den Mittelpunkt parallel zur Konstruktionsebene. Für einen
Kreis mit vorgegebenem \emph{Durchmesser} tippen Sie den halben Wert.

\subsection{Circle (3 points) -- Kreis durch drei Punkte}\label{sec:cmd-circle3p}
\cmd{circle3p}, Alias \code{c3}. Der Kreis durch drei beliebige Punkte im Raum; seine Ebene
ergibt sich aus den Punkten, nicht aus der Ansicht.
\begin{schritte}
\item \code{Circle: point 1 of 3}, \code{Circle: point 2 of 3}, \code{Circle: point 3 of 3}.
\end{schritte}
Liegen die drei Punkte auf einer Geraden oder fallen zwei zusammen, wird der dritte Punkt
abgelehnt.

\subsection{Arc -- Kreisbogen}\label{sec:cmd-arc}
\cmd{arc}, Alias \code{a}. Ein Kreisbogen durch Anfangspunkt, Endpunkt und einen dritten Punkt
auf dem Bogen.
\begin{schritte}
\item \code{Arc: start point}.
\item \code{Arc: end point}.
\item \code{Arc: point on the arc} -- ein beliebiger Punkt zwischen Anfang und Ende auf dem
  gewünschten Bogen; er legt Radius und Bogenseite fest.
\end{schritte}
Der Bogen kann beliebig im Raum liegen. Er läuft vom Anfangspunkt über den dritten Punkt zum
Endpunkt. Für einen Bogen mit bekanntem Radius $R$ und Pfeilhöhe berechnen Sie den dritten Punkt
(Beispiel in Abschnitt~\ref{sec:bsp-spiegel}).

\subsection{Helix -- Schraubenlinie}\label{sec:cmd-helix}
\cmd{helix}, Aliasse \code{h}, \code{he}. Eine zylindrische Schraubenlinie um eine beliebige
Achse; standardmäßig mit 5 Windungen.
\begin{schritte}
\item \code{Helix (5 turns): axis start (type a number to change the turns)} -- optional eine
  Windungszahl tippen (jede positive Zahl, auch nicht ganzzahlig), dann den Anfang der Achse
  setzen.
\item \code{Helix (5 turns): axis end} -- das Ende der Achse. Die Achslänge ist die Höhe der
  Helix; die Steigung ergibt sich als Höhe geteilt durch Windungen. Auch in diesem Schritt
  ändert eine Zahl noch die Windungszahl.
\item \code{Helix (5 turns): radius (point or number)} -- ein Punkt legt Radius und
  Startrichtung fest (sein Abstand von der Achse); eine getippte Zahl ist der Radius.
\end{schritte}
Eine Schraubenfeder modellieren Sie, indem Sie anschließend einen kleinen Kreis an den
Anfang der Helix setzen und ihn mit \gui{Sweep 1 rail} entlang der Helix führen.

\subsection{Ellipse}\label{sec:cmd-ellipse}
\cmd{ellipse}, Alias \code{el}. Eine Ellipse in der Konstruktionsebene, exakt (rationale
NURBS-Kurve wie der Kreis).
\begin{schritte}
\item \code{Ellipse: centre}.
\item \code{Ellipse: end of the first axis} -- legt Richtung und Länge der ersten Halbachse
  fest; eine getippte Zahl ist ihre Länge entlang der $u$-Achse der Ebene.
\item \code{Ellipse: second axis (point or length)} -- der Abstand des Punkts von der ersten
  Achse oder eine Zahl ist die zweite Halbachse.
\end{schritte}

\subsection{Spiral -- archimedische Spirale}\label{sec:cmd-spiral}
\cmd{spiral}, Alias \code{spi}. Eine ebene Spirale um einen Mittelpunkt in der
Konstruktionsebene, gegen den Uhrzeigersinn; standardmäßig 5 Windungen.
\begin{schritte}
\item \code{Spiral (5 turns): centre (type a number to change the turns)}.
\item \code{Spiral (5 turns): start radius (point or number)} -- ein Punkt legt auch die
  Startrichtung fest.
\item \code{Spiral (5 turns): end radius (point or number)} -- der Radius wächst (oder
  schrumpft) gleichmäßig bis hierher; 0 endet im Mittelpunkt.
\end{schritte}
Eine Spirale mit Steigung (konische Helix) zeichnen Sie weiterhin mit \gui{Helix}.

\subsection{Curve (control points) -- Kurve aus Kontrollpunkten}\label{sec:cmd-curve}
\cmd{curve}, Alias \code{crv}. Eine B-Spline-Kurve (Grad 3, mit weniger Punkten niedriger),
deren \emph{Kontrollpunkte} Sie setzen: Die Kurve beginnt und endet in den äußeren Punkten und
wird von den inneren nur angezogen -- sie verläuft glatter als eine interpolierende Kurve.
\begin{schritte}
\item \code{Curve: first control point}.
\item \code{Curve: next control point (Enter: done, snap to the first point: close)} -- die
  Vorschau zeigt Kurve und Kontrollpolygon. Ein Klick auf den ersten Punkt schließt die Kurve.
\item \keys{Enter} erzeugt die Kurve (ab zwei Punkten).
\end{schritte}

\subsection{Sketch -- Freihandkurve}\label{sec:cmd-sketch}
\cmd{sketch}, Alias \code{sk}. Klicken Sie den Anfang, führen Sie den Cursor die Kurve entlang
und beenden Sie mit einem Klick oder \keys{Enter}. Der Weg des Cursors wird durch eine kubische
Spline-Kurve angenähert (Abweichung höchstens 0,5\,\% seiner Länge); Anfang und Ende werden
exakt getroffen. Der Objektfang wirkt auch hier -- schalten Sie ihn für freies Zeichnen ab.

\subsection{Curve through points -- Kurve durch ausgewählte Punkte}\label{sec:cmd-crvpts}
\cmd{curve\_through\_points}, Aliasse \code{crvpts}, \code{curvethroughpoints}. Wählen Sie
zwei oder mehr Punktobjekte in der gewünschten Reihenfolge aus und starten Sie den Befehl: Er
legt sofort eine interpolierende Kurve durch sie (kubisch, mit weniger Punkten niedriger).

\subsection{Conic -- Kegelschnitt}\label{sec:cmd-conic}
\cmd{conic}. Ein Kegelschnittbogen von einem Anfangs- zu einem Endpunkt, tangential an zwei
gegebene Richtungen, exakt als rationale quadratische Kurve.
\begin{schritte}
\item \code{Conic: start point}, \code{Conic: end point}.
\item \code{Conic: apex (where the end tangents meet)} -- der Schnittpunkt der beiden
  Endtangenten.
\item \code{Conic: rho (0 … 1; 0.5: parabola) or a point it passes through} -- $\rho$ legt
  fest, wie weit die Kurve vom Sehnenmittelpunkt zum Apex reicht: unter 0,5 eine Ellipse, genau
  0,5 eine Parabel, darüber eine Hyperbel. Ein Punkt gibt $\rho$ durch seine Lage zwischen
  Sehnenmitte und Apex.
\end{schritte}

\subsection{Parabola -- Parabel}\label{sec:cmd-parabola}
\cmd{parabola}, Alias \code{para}. Ein zur Achse symmetrisches Parabelstück in der
Konstruktionsebene, exakt.
\begin{schritte}
\item \code{Parabola: vertex} -- der Scheitel.
\item \code{Parabola: focus (or focal length)} -- der Brennpunkt legt Achse und Brennweite $f$
  fest ($y = x^2/(4f)$); eine Zahl ist die Brennweite, geöffnet entlang der $v$-Achse.
\item \code{Parabola: end point (or half width)} -- der Abstand des Punkts von der Achse ist die
  halbe Breite; die Kurve endet auf seiner Seite.
\end{schritte}
Mit einer Parabel und \gui{Revolve} um ihre Achse entsteht ein Paraboloidspiegel.

\subsection{Hyperbola -- Hyperbel}\label{sec:cmd-hyperbola}
\cmd{hyperbola}, Alias \code{hyp}. Ein Ast der Hyperbel $x^2/a^2 - y^2/b^2 = 1$, symmetrisch zur
Achse, exakt.
\begin{schritte}
\item \code{Hyperbola: centre}.
\item \code{Hyperbola: vertex} -- der Scheitel des Astes: Achsrichtung und Halbachse $a$.
\item \code{Hyperbola: a point of the branch (farther from the centre than the vertex)} -- die
  Hyperbel geht durch diesen Punkt ($b$ ergibt sich daraus) und reicht symmetrisch bis zu
  seinem Abstand von der Achse.
\end{schritte}

\subsection{Interpolated curve -- interpolierende Spline-Kurve}\label{sec:cmd-spline}
\cmd{spline}, Aliasse \code{sp}, \code{interpcrv}. Eine glatte Freiformkurve, die exakt durch alle eingegebenen
Punkte verläuft (kubisch interpolierend; bei nur zwei bzw.\ drei Punkten Grad 1 bzw.\ 2).
\begin{schritte}
\item \code{Spline: first point}.
\item \code{Spline: next point (Enter: done)} -- beliebig viele weitere Punkte.
\item \keys{Enter} oder \maus{Rechtsklick} erzeugt die Kurve (ab zwei Punkten).
\end{schritte}
Splines eignen sich für Profile von Freiformflächen und als Leitkurven für Sweeps. Da sie die
Punkte interpoliert, schwingt die Kurve bei ungleichmäßig verteilten Punkten über; setzen Sie
Punkte möglichst gleichmäßig.

\subsection{Offset curve -- Parallelkurve}\label{sec:cmd-offset-curve}
\cmd{offset\_curve}, Alias \code{offset}. Erzeugt zu jeder ausgewählten Kurve eine
Parallelkurve in festem Abstand, in der Konstruktionsebene.
\begin{schritte}
\item Wählen Sie eine oder mehrere Kurven aus und starten Sie den Befehl.
\item \code{Offset curve (distance 1): click the side (type a number to change the distance)}
  -- tippen Sie zuerst den gewünschten Abstand (positive Zahl; Voreinstellung 1), dann klicken
  Sie auf die Seite, auf der die Parallelkurve entstehen soll. Die Vorschau zeigt sie bereits
  auf der Seite des Cursors.
\end{schritte}
Die Originale bleiben erhalten. Der Abstand wird senkrecht zur Kurve innerhalb der
Konstruktionsebene der Ansicht gemessen, in der Sie klicken -- zeichnen Sie ebene Kurven also
in der Ansicht, in deren Ebene sie liegen.

Kann die Parallelkurve nicht gebildet werden, nennt die Statusleiste den Grund, und der Befehl
wartet auf einen neuen Klick bzw.\ Abstand:
\begin{itemize}
\item \code{Offset curve: the distance must be positive (click the side)} -- Sie haben 0 oder
  einen negativen Abstand getippt; die Seite wählen Sie mit dem Klick.
\item \code{Offset curve: the curve does not lie in the construction plane of this view –
  click the side in a view parallel to the curve} -- klicken Sie in der Ansicht, in deren
  Ebene die Kurve liegt.
\item \code{Offset curve: the curve has corners – Explode it and offset the pieces, or fillet
  the corners first} -- Kurven mit Ecken (Polylinie, Rechteck, Polygon) zerlegen Sie mit
  \gui{Explode} oder runden ihre Ecken vorher mit \gui{Fillet curves} ab.
\item \code{Offset curve: distance <d> is larger than the radius of curvature on that side} --
  auf der Innenseite einer Krümmung darf der Abstand den Krümmungsradius nicht überschreiten;
  wählen Sie einen kleineren Abstand oder die andere Seite.
\end{itemize}

\subsection{Fillet curves -- Kurven verrunden}\label{sec:cmd-fillet-curve}
\cmd{fillet\_curve}, Alias \code{fillet}. Rundet Ecken mit einem Kreisbogen ab.
\begin{schritte}
\item Wählen Sie \emph{eine} oder \emph{zwei} ebene Kurven aus und starten Sie den Befehl.
\item \code{Fillet: radius (two curves: fillet between them; one curve: all its corners)} --
  tippen Sie den Radius.
\end{schritte}
\begin{itemize}
\item \textbf{Zwei Kurven} (etwa zwei Linien, die sich treffen oder schneiden): Ein
  Verrundungsbogen wird erzeugt, und beide Kurven werden bis zum Bogen gekürzt.
\item \textbf{Eine Kurve} (etwa eine Polylinie, ein Rechteck): \emph{Alle} Ecken der Kurve
  werden mit dem Radius verrundet; die Kurve wird ersetzt.
\end{itemize}
Passt der Radius nicht, nennt die Statusleiste den Grund, und der Befehl wartet auf einen
anderen Radius (oder \keys{Esc}):
\begin{itemize}
\item \code{Fillet: radius 5 too large for these curves} -- tippen Sie einen kleineren Radius.
\item \code{Fillet: the radius must be positive}.
\item \code{Fillet: the curves do not meet – extend them or pick two curves that cross}.
\item \code{Fillet: the curves are tangent where they meet (nothing to round)}.
\item \code{Fillet: the curve is not planar}.
\item \code{Fillet: the curve has no corners} -- die eine ausgewählte Kurve hat keine Ecke, die
  sich verrunden ließe.
\end{itemize}

% ============================================================================
\subsection{Chamfer curves -- Kurven fasen}\label{sec:cmd-chamfer}
\cmd{chamfer}, Aliasse \code{cha}, \code{chamfercorners}. Wie \gui{Fillet curves}, nur mit
einer geraden Fase statt eines Bogens.
\begin{schritte}
\item Wählen Sie eine oder zwei Kurven aus und starten Sie den Befehl.
\item \code{Chamfer: distance from the corner (two curves: between them; one curve: all its
  corners)} -- tippen Sie den Abstand der Fasenpunkte von der Ecke.
\end{schritte}
Bei zwei Kurven (die sich treffen oder schneiden) entsteht die Fasenlinie, und beide Kurven
werden bis zu ihr gekürzt; bei einer Kurve (Polylinie, Rechteck, Polygon) wird jede Ecke
gefast. Die Kurven müssen nicht eben sein. \code{filletcorners} startet entsprechend
\gui{Fillet curves}.

\subsection{Blend curves -- Übergangskurve}\label{sec:cmd-blend}
\cmd{blend}, Alias \code{blendcrv}. Eine glatte Verbindung zwischen den Enden zweier
ausgewählter Kurven, die beide Kurven ohne Knick fortsetzt.
\begin{schritte}
\item Wählen Sie zwei Kurven aus und starten Sie den Befehl. Die Vorschau zeigt die
  Übergangskurve zwischen den beiden einander nächsten Enden.
\item \code{Blend (G2): click near the end of the first curve, or Enter: the closest ends (type 1
  or 2: continuity)} -- \keys{Enter} verbindet die nächsten Enden; ein Klick wählt das Ende der
  ersten Kurve, ein zweiter Klick das der zweiten. Eine getippte 1 macht den Übergang
  tangentenstetig (G1, kubisch), 2 (Voreinstellung) auch krümmungsstetig (G2, Grad 5).
\end{schritte}
Die Originale bleiben erhalten; verbinden Sie sie bei Bedarf mit \gui{Join}.

\subsection{Tween curves -- Zwischenkurven}\label{sec:cmd-tween}
\cmd{tween}, Alias \code{tweencurves}. Kurven, die gleichmäßig von der ersten zur zweiten
ausgewählten Kurve überblenden.
\begin{schritte}
\item Wählen Sie zwei Kurven aus und starten Sie den Befehl.
\item \code{Tween: number of curves (Enter: 1)} -- tippen Sie die Anzahl oder drücken Sie
  \keys{Enter} für eine Kurve in der Mitte.
\end{schritte}
Läuft die zweite Kurve andersherum als die erste, wird sie vorher umgedreht.

\section{Kurven aus Objekten (Gruppe Curves from objects)}\label{sec:gruppe-derive}

Diese Befehle leiten Kurven aus vorhandenen Flächen, Körpern und Kurven ab. Sie arbeiten auf der
Auswahl; Flächen eines Körpers und Kanten können Sie auch einzeln als Teilobjekte wählen
(\keys{Strg+Shift}-Klick, Abschnitt~\ref{sec:teilobjekte}). Die Originale bleiben erhalten.

\subsection{Project curves, Pull curves -- Kurven auf Flächen abbilden}\label{sec:cmd-project}
\cmd{project} (Alias \code{proj}) und \cmd{pull}. Wählen Sie Kurven und die Flächen, auf die sie
sollen, und starten Sie den Befehl: Jede Kurve wird sofort auf jede Fläche abgebildet --
\gui{Project} entlang der Blickrichtung der aktiven Ansicht (Normale ihrer Konstruktionsebene),
\gui{Pull} senkrecht zur Fläche (nächste Punkte). Die neue Kurve liegt auf der ganzen Fläche; den
Beschnitt einer getrimmten Fläche berücksichtigt sie nicht. Trifft ein Teil der Kurve die Fläche
nicht, entsteht für dieses Paar keine Kurve.

\subsection{Section -- Schnitt}\label{sec:cmd-section}
\cmd{section}. Wählen Sie Flächen oder Körper und klicken Sie zwei Punkte einer Schnittlinie:
Geschnitten wird mit der Ebene durch diese Linie, senkrecht zur Ansicht. Ein geschlossener Körper
gibt geschlossene Schnittkurven, eine Fläche ihre Schnittkurven innerhalb ihres Beschnitts.

\subsection{Contour -- Höhenlinien}\label{sec:cmd-contour}
\cmd{contour}. Parallele Schnitte im festen Abstand über die ganze Auswahl.
\begin{schritte}
\item \code{Contour: base point} -- ein Punkt, durch den eine der Ebenen geht.
\item \code{Contour: direction of the contour planes' normal} -- die Richtung, in der die Ebenen
  aufeinander folgen.
\item \code{Contour: spacing (number)} -- der Abstand der Ebenen.
\end{schritte}
Höchstens 1000 Ebenen; liegt eine Ebene genau in einer ebenen Fläche eines Körpers, werden seine
Flächen einzeln geschnitten.

\subsection{Duplicate edge, Duplicate border -- Kanten und Ränder als Kurven}\label{sec:cmd-dupedge}
\cmd{dupedge}: Wählen Sie Kanten als Teilobjekte (\keys{Strg+Shift}-Klick) und starten Sie den
Befehl; jede wird eine eigenständige Kurve. \cmd{dupborder}: die Randkurven der ausgewählten
Flächen (der Außenrand und die Löcher, jeweils zu einer Kurve verbunden) bzw.\ die offenen
Kanten einer offenen Hülle.

\subsection{Extract isocurve, Extract wireframe -- Isokurven}\label{sec:cmd-iso}
\cmd{extract\_isocurve}, Aliasse \code{iso}, \code{extractisocurve}: Wählen Sie Flächen und
klicken Sie Punkte auf ihnen -- jeder Klick erzeugt die Isokurven durch den Punkt (Vorschau am
Cursor); eine getippte 1 nur die in $u$-Richtung konstante, 2 nur die in $v$, 3 beide.
\keys{Enter} beendet. \cmd{extract\_wireframe} (Alias \code{extractwireframe}) erzeugt sofort die
Isokurven an allen Knotenlinien der ausgewählten Flächen, die Ränder eingeschlossen. Beide
gelten für die ganze zugrunde liegende Fläche, auch bei getrimmten Flächen.

\subsection{Offset curve on surface -- Parallelkurve auf einer Fläche}\label{sec:cmd-offset-on-surface}
\cmd{offset\_on\_surface}, Aliasse \code{offsetcrvonsrf}, \code{offsetonsrf}. Wählen Sie eine
Kurve und die Fläche, auf der sie liegt, und tippen Sie den Abstand: Die neue Kurve hält ihn
\emph{auf der Fläche} (geodätisch), links der Kurve von der Flächennormalen aus gesehen; ein
negativer Abstand legt sie auf die andere Seite.

\subsection{Curve Boolean -- Boolesche Operationen mit Kurven}\label{sec:cmd-curve-boolean}
\cmd{curve\_boolean}, Aliasse \code{crvboolean}, \code{curveboolean}. Wählen Sie zwei oder mehr
geschlossene Kurven, die in der Konstruktionsebene der Ansicht liegen, und tippen Sie: 1
Vereinigung, 2 Schnittmenge, 3 Differenz (die zuerst gewählte minus die anderen). Die Ränder des
Ergebnisses entstehen als neue geschlossene Kurven; die Kurvenstücke bleiben exakt.

\section{Kurven bearbeiten (Gruppe Edit curves)}\label{sec:gruppe-curve-edit}

Diese Befehle ändern die ausgewählten Kurven an Ort und Stelle (ein Rückgängig-Schritt) oder
leiten Punkte von ihnen ab. Trim, Split, Join und Explode stehen in der Gruppe \gui{Edit}
(Abschnitt~\ref{sec:gruppe-edit}).

\subsection{Extend curve -- Kurve verlängern}\label{sec:cmd-extend}
\cmd{extend}, Alias \code{extendcrv}. \code{Extend: click near the end to extend}, dann
\code{Extend: length (number or a point)}: Die Kurve wird an diesem Ende um die Länge (bzw.\ den
Abstand des Punkts vom Ende) natürlich fortgesetzt -- ein Bogen bleibt ein Kreisbogen, eine
Gerade eine Gerade.

\subsection{Divide, Divide by length -- Punkte auf Kurven}\label{sec:cmd-divide}
\cmd{divide}: Punkte, die die ausgewählten Kurven in die getippte Anzahl gleich langer Stücke
teilen (bei geschlossenen Kurven ohne doppelten Anfangspunkt). \cmd{divide\_length} (Aliasse
\code{divlen}, \code{dividebylength}): Punkte im getippten Abstand entlang der Kurve vom Anfang
an.

\subsection{Rebuild curve, Change degree}\label{sec:cmd-rebuild}
\cmd{rebuild}: Die Kurve wird mit der getippten Anzahl Kontrollpunkte und dann dem Grad
(\keys{Enter}: 3) neu aufgebaut, mit gleichmäßigen Knoten; die Statusleiste nennt die größte
Abweichung. \cmd{change\_degree} (Alias \code{changedegree}): Ein höherer Grad ändert die Form
nicht; ein niedrigerer wird innerhalb von 0,1\,\% der Kurvengröße angenähert, sonst die Kurve mit
gleich vielen Kontrollpunkten neu aufgebaut.

\subsection{Flip direction, Close curve, Curve seam}\label{sec:cmd-flip}
\cmd{flip} (Alias \code{dir}) kehrt die Richtung der ausgewählten Kurven und die Normalen
ausgewählter Flächen um. \cmd{close} (Alias \code{closecrv}) schließt offene Kurven: Liegen die
Enden nahe beieinander (höchstens 1\,\% der Länge), wird das Ende auf den Anfang gezogen, sonst
ein gerades Stück angefügt. \cmd{seam} (Alias \code{crvseam}) verlegt Anfang und Ende einer
geschlossenen Kurve an den geklickten Punkt.

\subsection{Insert knot, Remove knot -- Knoten}\label{sec:cmd-knots}
\cmd{insert\_knot} (Alias \code{insertknot}): Jeder Klick auf die ausgewählte Kurve fügt dort
einen Knoten ein -- die Form bleibt, die Kurve bekommt einen Kontrollpunkt mehr.
\cmd{remove\_knot} (Alias \code{removeknot}): Jeder Klick entfernt den inneren Knoten, dessen
Kurvenpunkt am nächsten liegt; dabei ändert sich die Form. \keys{Enter} beendet beide.

\subsection{Smooth curve, Fair curve -- glätten}\label{sec:cmd-smooth}
\cmd{smooth}: Ein Glättungsschritt -- jeder innere Kontrollpunkt rückt um den getippten Faktor
(0 bis 1, \keys{Enter}: 0,5) zum Mittel seiner Nachbarn; die Enden bleiben. Mehrfach anwenden
glättet stärker. \cmd{fair}: Glättet schrittweise, solange die Kurve innerhalb der getippten
Toleranz von der ursprünglichen bleibt -- Wellen verschwinden, die Form bleibt.

\section{Flächen (Gruppe Surfaces)}\label{sec:gruppe-surfaces}
\index{Fläche}

Bis auf Plane erzeugen alle Befehle dieser Gruppe Flächen aus vorhandenen Kurven oder Flächen,
die Sie vorher auswählen. Das Ergebnis sind \emph{offene Flächen}, keine Volumenkörper; mit
\gui{Cap planar holes} und \gui{Join} machen Sie daraus geschlossene Körper
(Abschnitt~\ref{sec:cmd-cap}). Die Ausgangskurven bleiben jeweils erhalten.

\subsection{Plane -- ebene Rechteckfläche}\label{sec:cmd-plane}
\cmd{plane}. Eine rechteckige ebene Fläche in der Konstruktionsebene, mit den Kanten parallel zu
deren Achsen.
\begin{schritte}
\item \code{Plane: first corner}.
\item \code{Plane: opposite corner}.
\end{schritte}
Die Bedienung entspricht Rectangle, nur entsteht eine Fläche statt einer Randkurve. Ebene
Flächen sind nützlich als Schirme, Detektorflächen oder Schnittwerkzeuge für Split und Trim.

\subsection{Extrude -- Extrusion}\label{sec:cmd-extrude}
\cmd{extrude}, Alias \code{ext}. Zieht jede ausgewählte Kurve geradlinig zu einer Fläche aus.
\begin{schritte}
\item Wählen Sie eine oder mehrere Kurven aus und starten Sie den Befehl.
\item \code{Extrude: distance along the construction-plane normal, or start of the direction}
  -- tippen Sie eine Distanz, \emph{oder} setzen Sie den Anfangspunkt einer Richtung.
\item Falls Sie einen Punkt gesetzt haben: \code{Extrude: end of the direction} -- der
  Endpunkt. Der Vektor vom ersten zum zweiten Punkt ist Richtung und Länge der Extrusion. Statt
  des Endpunkts können Sie auch eine Länge tippen; sie wirkt dann in Richtung des Cursors.
\end{schritte}
Eine \emph{eingetippte} Distanz wirkt entlang der Normalen $n$ der Konstruktionsebene der
Ansicht, in der Sie arbeiten -- der aktiven Ansicht beim Start des Befehls, danach der Ansicht,
über der die Maus zuletzt war (Abschnitt~\ref{sec:mod-konstruktionsebene}):

\begin{referenz}{lX}
\toprule
Ansicht & Richtung einer positiven Distanz \\
\midrule\endhead
Top, Perspective & $+z$ \\
Front & $-y$ (zum Betrachter) \\
Left & $-x$ (zum Betrachter) \\
\bottomrule
\end{referenz}

Eine negative Distanz extrudiert in die Gegenrichtung. Die Distanz 0 wird mit
\code{Extrude: the distance must not be 0} abgelehnt. Für eine beliebige andere Richtung geben
Sie zwei Punkte ein, etwa \code{0,0,0} und \code{0,5,10} für eine schräge Extrusion.
Eine geschlossene Kurve (Kreis, Rechteck, Polygon, geschlossene Polylinie) ergibt einen offenen
Mantel; mit \gui{Cap planar holes} wird daraus ein Volumenkörper.

\subsection{Revolve -- Rotationsfläche}\label{sec:cmd-revolve}
\cmd{revolve}, Alias \code{rev}. Dreht jede ausgewählte Kurve um eine Achse. Das ist der
wichtigste Befehl für rotationssymmetrische Optik: Linsen- und Spiegelflächen, Fassungen,
Reflektoren.
\begin{schritte}
\item Wählen Sie das Profil (eine oder mehrere Kurven) aus und starten Sie den Befehl.
\item \code{Revolve: start of the axis} -- erster Punkt der Drehachse.
\item \code{Revolve: end of the axis} -- zweiter Punkt der Drehachse. Die Vorschau zeigt das
  Profil in \SI{90}{\degree}-Schritten um die Achse.
\item \code{Revolve: angle in degrees (Enter: 360)} -- tippen Sie einen Winkel (nicht 0,
  höchstens \SI{360}{\degree} im Betrag) oder drücken Sie \keys{Enter} für eine volle Umdrehung.
\end{schritte}
Fällt das Achsende auf den Achsanfang, meldet die Statusleiste \code{Revolve: the end of the
axis must differ from its start}; ein Winkel außerhalb des Bereichs ergibt \code{Revolve: the
angle must be between -360 and 360 degrees and not 0}, ein weiterer Punkt statt des Winkels
\code{Revolve: type the angle in degrees, or Enter for 360}. Der Befehl wartet jeweils auf eine
passende Eingabe.
Positive Winkel drehen gegen den Uhrzeigersinn um die Achsrichtung (vom Anfangs- zum Endpunkt,
Rechte-Hand-Regel), negative Winkel in die andere Richtung. Für die optische Achse $z$ wählen
Sie z.\,B.\ \code{0,0,0} und \code{0,0,1} als Achspunkte. Jede Profilkurve ergibt eine eigene
Fläche; verbinden Sie die Profilkurven vorher mit \gui{Join}, wenn Sie eine einzige Fläche
wünschen.

\subsection{Loft -- Fläche durch Querschnitte}\label{sec:cmd-loft}
\cmd{loft}. Eine glatte Fläche durch zwei oder mehr Querschnittskurven.
\begin{schritte}
\item Wählen Sie die Querschnitte \emph{in der gewünschten Reihenfolge} aus (Klick,
  \keys{Shift}+Klick, \ldots).
\item Starten Sie \gui{Loft}. Die Fläche entsteht sofort.
\end{schritte}
Die Querschnitte sollten gleich orientiert sein (alle im selben Drehsinn bzw.\ in derselben
Richtung laufen), sonst verdreht sich die Fläche. Mit weniger als zwei Kurven erscheint
\code{Loft: select two or more section curves (in order)}.

\subsection{Sweep 1 rail -- Profil entlang einer Leitkurve}\label{sec:cmd-sweep1}
\cmd{sweep1}. Führt ein Profil entlang einer Leitkurve (Rail).
\begin{schritte}
\item Wählen Sie genau zwei Kurven aus: zuerst das \emph{Profil}, dann die \emph{Leitkurve}.
\item Starten Sie \gui{Sweep 1 rail}. Die Fläche entsteht sofort.
\end{schritte}
Bei falscher Anzahl: \code{Sweep 1 rail: select the profile, then the rail}. Legen Sie das
Profil an den Anfang der Leitkurve, etwa einen Kreis senkrecht zur Helix für eine Feder.

\subsection{Sweep 2 rails -- Profil entlang zweier Leitkurven}\label{sec:cmd-sweep2}
\cmd{sweep2}. Führt ein Profil entlang zweier Leitkurven, die zugleich die Breite steuern.
\begin{schritte}
\item Wählen Sie genau drei Kurven aus: zuerst das \emph{Profil}, dann \emph{Leitkurve 1}, dann
  \emph{Leitkurve 2}.
\item Starten Sie \gui{Sweep 2 rails}.
\end{schritte}
Das Profil sollte die beiden Leitkurven an ihren Anfängen berühren.

\subsection{Coons patch -- Fläche aus vier Randkurven}\label{sec:cmd-coons}
\cmd{coons}. Eine Fläche, die genau von vier Randkurven begrenzt wird.
\begin{schritte}
\item Wählen Sie genau vier Kurven aus, die ein geschlossenes Viereck bilden: zwei
  gegenüberliegende, dann die beiden anderen.
\item Starten Sie \gui{Coons patch}.
\end{schritte}
Die Richtung der einzelnen Kurven wird automatisch angepasst. Die Ecken müssen sich treffen
(fangen Sie die Endpunkte beim Zeichnen mit \gui{End}).

\subsection{Gordon surface -- Fläche aus einem Kurvennetz}\label{sec:cmd-gordon}
\cmd{gordon}, Alias \code{network}. Eine Fläche durch ein Netz sich kreuzender Kurven
(mindestens $2 \times 2$).
\begin{schritte}
\item Wählen Sie alle Kurven des Netzes aus, in beliebiger Reihenfolge (mindestens vier).
\item Starten Sie \gui{Gordon surface}.
\end{schritte}
Das Programm sortiert die Kurven selbst in zwei Scharen: Die Kurven, die die zuerst gewählte
Kurve schneiden, bilden die eine Schar, die übrigen die andere. Jede Kurve der einen Schar muss
jede Kurve der anderen schneiden. Bilden die Kurven kein Netz, erscheint
\code{Gordon surface: the curves do not form a network}.

\subsection{Offset surface -- Parallelfläche}\label{sec:cmd-offset-surface}
\cmd{offset\_surface}, Alias \code{offsetsrf}. Erzeugt zu jeder ausgewählten Fläche eine
Parallelfläche im festen Abstand entlang der Flächennormalen.
\begin{schritte}
\item Wählen Sie eine oder mehrere \emph{unbeschnittene} Flächen aus und starten Sie den
  Befehl.
\item \code{Offset surface: distance along the normal (negative: other side)} -- tippen Sie den
  Abstand; ein negatives Vorzeichen legt die Parallelfläche auf die andere Seite.
\end{schritte}
\begin{hinweis}
\enquote{Unbeschnitten} sind Flächen, die direkt aus Plane, Extrude, Revolve, Loft, Sweep,
Coons, Gordon oder Offset entstanden sind. Flächen eines Volumenkörpers (Box, Kugel, Ergebnis
einer Booleschen Operation oder von Join/Cap) und beschnittene Flächen werden nicht
angenommen; dann erscheint \code{Offset surface: select (untrimmed) surfaces}.
\end{hinweis}

\subsection{Fillet surfaces -- Flächen verrunden}\label{sec:cmd-fillet-surface}
\cmd{fillet\_surface}, Alias \code{filletsrf}. Erzeugt eine Rundungsfläche (Rollkugel-Fillet)
zwischen zwei Flächen.
\begin{schritte}
\item Wählen Sie genau zwei unbeschnittene Flächen aus und starten Sie den Befehl.
\item \code{Fillet surfaces: radius} -- tippen Sie den Radius.
\end{schritte}
Auf welcher Seite der Flächen die Rundung liegt, ermittelt das Programm selbst (es probiert die
möglichen Seiten durch). Die Ausgangsflächen werden nicht beschnitten; das erledigen Sie bei
Bedarf mit Trim.

\subsection{Cap planar holes -- ebene Öffnungen schließen}\label{sec:cmd-cap}
\cmd{cap}. Schließt alle ebenen offenen Ränder der ausgewählten Flächen mit beschnittenen
ebenen Deckeln und näht alles zu \emph{einem} Körper zusammen. Typisch: Eine extrudierte
geschlossene Kurve wird so zum Volumenkörper.
\begin{schritte}
\item Wählen Sie die Fläche(n) aus und starten Sie \gui{Cap planar holes} (Gruppe Surfaces).
\end{schritte}
Das Ergebnis ersetzt das zuerst ausgewählte Objekt (Name, Layer und Farbe bleiben), die übrigen
ausgewählten Objekte gehen darin auf. Die Statusleiste meldet
\code{Capped: closed solid, volume ...}, wenn ein geschlossener Körper entstanden ist, sonst
\code{Capped (still open)}. Gibt es keinen ebenen offenen Rand, erscheint
\code{Cap planar holes: no planar open loop to cap}.

% ============================================================================
\section{Volumenkörper (Gruppe Solids)}\label{sec:gruppe-solids}
\index{Volumenkörper}

Die Grundkörper werden als geschlossene Hüllen (B-Rep-Schalen) mit exakten Flächen erzeugt.
Box, Cylinder, Cone und Torus beziehen sich auf die Konstruktionsebene der Ansicht, in der Sie
den ersten Punkt setzen; die Höhe zählt positiv entlang der Normalen $n$
(Abschnitt~\ref{sec:mod-konstruktionsebene}), negative Höhen bauen den Körper auf die andere Seite.

\subsection{Box -- Quader}\label{sec:cmd-box}
\cmd{box}, Alias \code{b}.
\begin{schritte}
\item \code{Box: first corner of the base}.
\item \code{Box: opposite corner of the base} -- die Grundfläche ist ein Rechteck in der
  Konstruktionsebene; es muss in beiden Richtungen eine Ausdehnung haben.
\item \code{Box: height (point or number)} -- tippen Sie die Höhe (Vorzeichen = Richtung) oder
  klicken Sie einen Punkt; dann zählt dessen Abstand von der Grundfläche entlang der Normalen.
\end{schritte}
Die Kanten des Quaders liegen parallel zu den Weltachsen.

\subsection{Sphere -- Kugel}\label{sec:cmd-sphere}
\cmd{sphere}, Alias \code{sph}.
\begin{schritte}
\item \code{Sphere: centre}.
\item \code{Sphere: radius (point or number)} -- Radius tippen oder einen Punkt auf der Kugel
  setzen (Abstand im Raum).
\end{schritte}

\subsection{Cylinder -- Zylinder}\label{sec:cmd-cylinder}
\cmd{cylinder}, Alias \code{cyl}. Ein Kreiszylinder, dessen Achse senkrecht zur
Konstruktionsebene steht.
\begin{schritte}
\item \code{Cylinder: centre of the base} -- Mittelpunkt der Grundfläche.
\item \code{Cylinder: radius (point or number)}.
\item \code{Cylinder: height (point or number)} -- mit Vorzeichen.
\end{schritte}

\subsection{Cone -- Kegel}\label{sec:cmd-cone}
\cmd{cone}. Ein spitzer Kreiskegel; Eingabe wie Cylinder:
\code{Cone: centre of the base}, \code{Cone: radius (point or number)},
\code{Cone: height (point or number)}. Die Spitze liegt auf der Achse in der angegebenen Höhe.

\subsection{Torus -- Ring}\label{sec:cmd-torus}
\cmd{torus}, Alias \code{tor}. Ein Torus, dessen Mittelebene in der Konstruktionsebene liegt.
\begin{schritte}
\item \code{Torus: centre}.
\item \code{Torus: major radius (point or number)} -- Radius des Mittelkreises.
\item \code{Torus: tube radius (point or number)} -- Radius des Rohrquerschnitts. Ein Punkt
  bestimmt ihn als Abstand vom Mittelkreis. Er muss kleiner als der große Radius sein.
\end{schritte}

\subsection{Boolesche Operationen}\label{sec:boolesche}
\index{Boolesche Operation}

Die drei Booleschen Operationen verknüpfen geschlossene Volumenkörper zu einem neuen Körper.
Sie arbeiten sofort auf der Auswahl, ohne weitere Eingabe.

\begin{referenz}{lllX}
\toprule
Werkzeug & Name & Alias & Ergebnis \\
\midrule\endhead
\gui{Boolean union} & \cmd{union} & \code{bu} & Vereinigung aller ausgewählten Körper. \\
\gui{Boolean difference} & \cmd{difference} & \code{bd} & Der \emph{zuerst} ausgewählte Körper
  minus alle anderen ausgewählten Körper. \\
\gui{Boolean intersection} & \cmd{intersection} & \code{bi} & Der Bereich, der in
  \emph{allen} ausgewählten Körpern liegt. \\
\bottomrule
\end{referenz}

\begin{schritte}
\item Wählen Sie mindestens zwei Volumenkörper aus. Für die Differenz zuerst den Körper, aus dem
  ausgeschnitten wird, dann mit \keys{Shift}+Klick die Werkzeugkörper.
\item Starten Sie die Operation. Das Ergebnis ersetzt die Eingangskörper.
\end{schritte}

Der neue Körper übernimmt Name, Layer und Farbe des zuerst ausgewählten Körpers; die
Eingangskörper werden entfernt. Mit \keys{Strg+Z} erhalten Sie sie zurück. Bei mehr als zwei
Körpern werden Vereinigung und Schnitt nacheinander über alle Körper gebildet.

Mögliche Meldungen:
\begin{itemize}
\item \code{Select at least two solids} bzw.\ \code{Select the solid to cut first, then the
  solids to subtract} -- es sind weniger als zwei Körper ausgewählt (Kurven und Flächen zählen
  nicht).
\item \code{Boolean ... failed: Booleans need closed solids} -- einer der Körper ist eine
  offene Hülle, etwa nach Join aus Flächen, die keinen geschlossenen Körper bilden. Einzelne
  Flächen wie eine Extrusion ohne Deckel zählen gar nicht als Körper; schließen Sie sie vorher
  mit Cap.
\item \code{Boolean ... failed: the result is empty} -- die Körper überlappen sich nicht (bei
  Schnitt) oder der Werkzeugkörper umschließt den Grundkörper ganz (bei Differenz).
\end{itemize}

\begin{tipp}
Vermeiden Sie bei Differenzen genau aufeinanderliegende Flächen: Lassen Sie einen
Bohrungszylinder ein Stück über den Grundkörper hinausragen, statt ihn exakt bündig enden zu
lassen (Beispiel in Abschnitt~\ref{sec:bsp-platte}).
\end{tipp}

% ============================================================================
\section{Transformieren (Gruppe Transform)}\label{sec:gruppe-transform}
\index{Transformation}

Alle Befehle dieser Gruppe wirken auf die ausgewählten Objekte -- Kurven, Flächen, Körper,
Punkte (Abschnitt~\ref{sec:auswahlbefehle}). Die Transformationen sind exakt: NURBS ändern
unter diesen Abbildungen ihre Form nicht. Während der Eingabe zeigt die Vorschau die
transformierten Objekte als Drahtmodell. Move, Rotate und Scale verändern die Objekte selbst;
Copy, Mirror und die Reihen (Arrays) erzeugen Kopien, die Name, Layer und Farbe des Originals
übernehmen.

\subsection{Move -- Verschieben}\label{sec:cmd-move}
\cmd{move}, Alias \code{m}.
\begin{schritte}
\item \code{Move: point to move from} -- Bezugspunkt, z.\,B.\ eine gefangene Ecke.
\item \code{Move: point to move to} -- Zielpunkt. Die Objekte werden um den Vektor zwischen
  beiden Punkten verschoben.
\end{schritte}
Für eine Verschiebung um einen bekannten Vektor tippen Sie als zweiten Punkt \code{@dx,dy,dz};
für eine Verschiebung um eine Länge in Cursorrichtung tippen Sie nur die Länge.

\subsection{Copy -- Kopieren}\label{sec:cmd-copy}
\cmd{copy}, Aliasse \code{co}, \code{cp}. Wie Move, aber die Originale bleiben, und Sie können
mehrere Kopien nacheinander ablegen.
\begin{schritte}
\item \code{Copy: point to move from}.
\item \code{Copy: point to copy to (Enter: done)} -- jeder Punkt legt eine Kopie ab, jeweils
  bezogen auf den ersten Punkt. Beenden mit \keys{Enter}.
\end{schritte}

\subsection{Rotate -- Drehen}\label{sec:cmd-rotate}
\cmd{rotate}, Alias \code{ro}. Dreht um eine Achse durch einen Punkt, parallel zur Normalen der
Konstruktionsebene (in Top also um eine $z$-parallele Achse).
\begin{schritte}
\item \code{Rotate: centre (axis: normal of the construction plane)} -- Drehzentrum.
\item \code{Rotate: angle in degrees or first reference point} -- tippen Sie den Winkel (in
  Grad, positiv gegen den Uhrzeigersinn um die Normale gesehen, also so, wie es in der Ansicht
  erscheint) \emph{oder} setzen Sie einen ersten Bezugspunkt.
\item Nur nach einem Bezugspunkt: \code{Rotate: second reference point} -- die Objekte werden
  um den Winkel zwischen den Richtungen Zentrum--Bezugspunkt~1 und Zentrum--Bezugspunkt~2
  gedreht. So richten Sie ein Objekt an vorhandener Geometrie aus.
\end{schritte}
Um eine andere Achsrichtung zu verwenden, arbeiten Sie in der passenden Ansicht: in Front um
eine $y$-parallele, in Left um eine $x$-parallele Achse.

\subsection{Scale -- Skalieren}\label{sec:cmd-scale}
\cmd{scale}, Alias \code{sc}. Gleichmäßige Skalierung in allen drei Richtungen um einen
Basispunkt.
\begin{schritte}
\item \code{Scale: base point} -- der Punkt, der fest bleibt.
\item \code{Scale: factor or first reference point} -- tippen Sie den Faktor (größer 0;
  \code{2} verdoppelt, \code{0.5} halbiert) \emph{oder} setzen Sie einen Bezugspunkt.
\item Nur nach einem Bezugspunkt: \code{Scale: second reference point} -- der Faktor ist das
  Verhältnis der Abstände Basispunkt--Punkt~2 zu Basispunkt--Punkt~1.
\end{schritte}
Mit den Bezugspunkten skalieren Sie ein Objekt auf ein Zielmaß, ohne den Faktor auszurechnen:
Basispunkt und Bezugspunkt~1 an die Enden einer Kante fangen, Bezugspunkt~2 im gewünschten
Abstand tippen (etwa \code{@50,0}).

\subsection{Mirror -- Spiegeln}\label{sec:cmd-mirror}
\cmd{mirror}, Alias \code{mi}. Erzeugt gespiegelte \emph{Kopien}; die Originale bleiben.
\begin{schritte}
\item \code{Mirror: start of the mirror line (copies, plane through the line and the
  construction-plane normal)}.
\item \code{Mirror: end of the mirror line}.
\end{schritte}
Gespiegelt wird an der Ebene, die die Spiegellinie enthält und senkrecht zur
Konstruktionsebene steht. In Top spiegelt eine Linie entlang der $y$-Achse also an der
$yz$-Ebene. Die Spiegellinie darf nicht parallel zur Normalen liegen.

\subsection{Linear array -- lineare Reihe}\label{sec:cmd-array-linear}
\cmd{array\_linear}, Alias \code{array}. Kopien in gleichen Schritten entlang einer Richtung.
\begin{schritte}
\item \code{Linear array (3 items): base point (type a number to change the count)} --
  optional die Gesamtzahl tippen (ganze Zahl ab 2, \emph{einschließlich} des Originals), dann
  den Basispunkt setzen.
\item \code{Linear array (3 items): step to the next item} -- der Punkt, an dem die erste
  Kopie liegen soll; der Vektor Basispunkt--Punkt ist der Schritt.
\end{schritte}
Mit \code{@10,0} als zweitem Punkt entsteht eine Reihe im Abstand 10 in $u$-Richtung.

\subsection{Grid array -- rechteckige Matrix}\label{sec:cmd-array-grid}
\cmd{array\_grid}. Kopien in einem Raster aus zwei Richtungen; standardmäßig $3 \times 3$.
\begin{schritte}
\item \code{Grid array (3 $\times$ 3): base point (type two numbers to change the counts)} --
  optional zwei Zahlen nacheinander tippen (jeweils mit \keys{Enter}): erst die Anzahl in
  Richtung~1, dann in Richtung~2 (ganze Zahlen ab 1). Dann den Basispunkt setzen.
\item \code{Grid array (3 $\times$ 3): step in the first direction} -- Schritt in Richtung~1.
\item \code{Grid array (3 $\times$ 3): step in the second direction} -- Schritt in
  Richtung~2.
\end{schritte}
Beide Schritte beziehen sich auf den Basispunkt, so dass \code{@30,0} und \code{@0,30} ein
quadratisches Raster mit 30 Einheiten Abstand ergeben.

\subsection{Polar array -- Kreisreihe}\label{sec:cmd-array-polar}
\cmd{array\_polar}, Alias \code{arraypolar}. Kopien gleichmäßig auf einem Kreisbogen um eine
Achse parallel zur Normalen der Konstruktionsebene; standardmäßig 6 Elemente.
\begin{schritte}
\item \code{Polar array (6 items): centre (type a number to change the count; axis:
  construction-plane normal)} -- optional die Gesamtzahl (ganze Zahl ab 2, einschließlich
  Original), dann das Zentrum setzen.
\item \code{Polar array (6 items): total angle in degrees (Enter: 360)} -- tippen Sie den
  Gesamtwinkel (nicht 0, höchstens \SI{360}{\degree} im Betrag) oder drücken Sie \keys{Enter}
  für einen Vollkreis.
\end{schritte}
Beim Vollkreis fällt die letzte Kopie nicht auf das Original: 6 Elemente liegen im Abstand von
\SI{60}{\degree}. Die Vorschau zeigt immer die Verteilung auf dem Vollkreis.

% ============================================================================
\section{Bearbeiten (Gruppe Edit)}\label{sec:gruppe-edit}

\subsection{Trim -- Stutzen}\label{sec:cmd-trim}
\cmd{trim}, Alias \code{tr}. Schneidet ein Objekt an anderen Objekten und entfernt den
angeklickten Teil.
\begin{schritte}
\item Wählen Sie zuerst das zu stutzende Objekt aus, dann mit \keys{Shift}+Klick die
  Schneidobjekte.
\item \code{Trim: click the part of the first selected object to remove} -- klicken Sie auf das
  Stück, das wegfallen soll.
\end{schritte}
Zulässige Kombinationen:
\begin{itemize}
\item \textbf{Kurve} an Kurven und/oder Flächen: Die Kurve wird an allen Schnittpunkten
  geteilt; das angeklickte Stück wird entfernt, die übrigen Stücke bleiben als einzelne Kurven.
\item \textbf{Fläche} (eine einzelne Fläche) an einer Schnittfläche: Die Fläche wird entlang der
  Schnittkurve geteilt; der angeklickte Teil wird entfernt.
\end{itemize}
Schneiden die Schneidobjekte die Kurve nicht, wird der Klick mit \code{Trim: the cutting
objects do not cross the curve} abgelehnt; liegt der Klick bei einer Fläche auf keinem Teil, mit
\code{Trim: click on the part to remove}. Klicken Sie erneut, oder brechen Sie mit \keys{Esc} ab
und prüfen Sie die Auswahl.

\subsection{Split -- Teilen}\label{sec:cmd-split}
\cmd{split}. Teilt ein Objekt, ohne etwas zu entfernen.
\begin{itemize}
\item \textbf{Eine Kurve an einem Punkt:} Wählen Sie genau eine Kurve aus und starten Sie
  Split. Aufforderung \code{Split: click the point on the curve} -- klicken Sie (am besten mit
  dem Fang \gui{Near}, \gui{Mid} oder \gui{Int}) auf die Teilungsstelle. Es entstehen zwei
  Kurven. Ein Punkt an einem Kurvenende wird mit \code{Split: the point is at an end of the
  curve} abgelehnt; der Befehl wartet auf einen anderen Punkt.
\item \textbf{Eine Fläche an einer Schnittfläche:} Wählen Sie zuerst die zu teilende Fläche,
  dann die Schnittfläche, und starten Sie Split. Die Fläche wird sofort entlang der Schnittkurve
  geteilt. Schneiden sich die Flächen nicht: \code{Split: the surfaces do not intersect}.
\end{itemize}
Die Teile übernehmen Name, Layer und Farbe des Originals.

\subsection{Join -- Verbinden}\label{sec:cmd-join}
\cmd{join}, Alias \code{j}. Verbindet mehrere Objekte zu einem.
\begin{itemize}
\item \textbf{Kurven:} Die ausgewählten Kurven werden Ende an Anfang zu \emph{einer} Kurve
  verkettet -- in beliebiger Auswahlreihenfolge und Laufrichtung; Kurven werden bei Bedarf
  umgedreht. Die Enden müssen sich berühren, sonst erscheint
  \code{Join: the curves do not form one chain}.
\item \textbf{Flächen:} Flächen und Flächen von Körpern werden entlang gemeinsamer Kanten zu
  einer Hülle vernäht. Die Statusleiste meldet die Anzahl der Flächen und ob ein geschlossener
  Körper entstanden ist, z.\,B.\ \code{Joined 6 faces, closed solid}.
\end{itemize}
Kurven und Flächen lassen sich nicht gemischt verbinden (\code{select either curves or
surfaces}); nötig sind mindestens zwei Objekte. Das Ergebnis ersetzt das zuerst ausgewählte
Objekt; die anderen gehen darin auf.

\subsection{Explode -- Zerlegen}\label{sec:cmd-explode}
\cmd{explode}, Alias \code{x}. Das Gegenstück zu Join.
\begin{itemize}
\item \textbf{Körper und Hüllen} werden in ihre einzelnen Flächen zerlegt.
\item \textbf{Kurven} werden an ihren Knicken zerlegt: eine Polylinie, ein Rechteck oder Polygon
  in einzelne Strecken. Auch Kreise und Bögen über \SI{90}{\degree} sind intern aus
  Viertelbögen zusammengesetzt und zerfallen dabei in diese Teilbögen.
\end{itemize}
Die Statusleiste meldet die Anzahl der Teile. Glatte Kurven ohne Knick bleiben unverändert;
gibt es gar nichts zu zerlegen, erscheint \code{Explode: nothing to explode}.

\subsection{Delete -- Löschen}\label{sec:cmd-delete}
Das Werkzeug \gui{Delete} löscht die ausgewählten Objekte -- wie \menu{Edit,Delete},
\keys{Entf} oder die Rücktaste. Es ist kein Befehl der Befehlszeile. Einzelne Flächen, Kanten
oder Ecken eines Körpers lassen sich nicht für sich löschen; dafür zerlegen Sie den Körper
zuerst mit Explode.

% ============================================================================
\section{Analysieren (Gruppe Analyze)}\label{sec:gruppe-analyze}
\index{Messen}

Die Analysebefehle verändern das Modell nicht (mit Ausnahme von Intersect, das Punkte und
Kurven erzeugt). Ihre Ergebnisse stehen in der Statusleiste. Zahlen werden mit bis zu acht
(Distance: zehn) gültigen Stellen ausgegeben.

\subsection{Distance -- Abstand}\label{sec:cmd-distance}
\cmd{distance}, Aliasse \code{dist}, \code{di}.
\begin{schritte}
\item \code{Distance: first point}.
\item \code{Distance: second point}.
\end{schritte}
Ergebnis: \code{Distance 25   (dx 15, dy 20, dz 0)} -- der räumliche Abstand und seine
Komponenten. Die Meldung bleibt in der Statusleiste stehen, bis eine neue erscheint. Fangen Sie
die Punkte, um exakte Abstände zu erhalten.

\subsection{Angle -- Winkel}\label{sec:cmd-angle}
\cmd{angle}, Alias \code{ang}.
\begin{schritte}
\item \code{Angle: vertex} -- der Scheitelpunkt.
\item \code{Angle: point on the first leg}.
\item \code{Angle: point on the second leg}.
\end{schritte}
Ergebnis: \code{Angle 90}\si{\degree} -- der Winkel zwischen den Schenkeln im Raum (0 bis
\SI{180}{\degree}).

\subsection{Curvature -- Krümmung}\label{sec:cmd-curvature}
\cmd{curvature}, Alias \code{curv}.
\begin{schritte}
\item Wählen Sie genau eine Kurve oder eine Fläche (bzw.\ einen Körper) aus und starten Sie den
  Befehl.
\item \code{Curvature: point on the selected curve or surface} -- klicken Sie auf die Stelle,
  an der gemessen werden soll. Ausgewertet wird der nächstgelegene Punkt der Kurve bzw.\
  Fläche.
\end{schritte}
Ergebnis bei einer Kurve: \code{Curvature 0.02   radius 50} (Krümmung und Krümmungsradius) bzw.\
\code{Curvature 0 (straight)}. Bei einer Fläche: \code{k1 ...   k2 ...   Gaussian ...   mean
...} -- die beiden Hauptkrümmungen, die Gaußsche und die mittlere Krümmung. Bei einem Körper wird
die Fläche verwendet, die dem Klickpunkt am nächsten liegt. So prüfen Sie zum Beispiel den
Scheitelradius einer modellierten Linsenfläche.

\subsection{Area -- Fläche und Länge}\label{sec:cmd-area}
\cmd{area}. Summiert die Flächeninhalte aller ausgewählten Flächen und Körper und die Längen
aller ausgewählten Kurven. Wählen Sie die Objekte aus und starten Sie den Befehl. Ergebnis:
\code{Area ...   Length ...} (nur die Teile, die vorkommen).

\subsection{Volume / mass properties -- Volumen und Schwerpunkt}\label{sec:cmd-volume}
\cmd{mass\_properties}, Alias \code{volume}. Berechnet das Gesamtvolumen und den gemeinsamen
Schwerpunkt der ausgewählten \emph{geschlossenen} Körper. Ergebnis:
\code{Volume ...   centroid (x, y, z)}. Offene Flächen werden übergangen; ist kein geschlossener
Körper ausgewählt, erscheint \code{Volume / mass properties: select closed solids}.
Volume eignet sich gut, um zu prüfen, ob ein Körper wirklich geschlossen ist.

\subsection{Intersection check -- Schnittpunkte und Schnittkurven}\label{sec:cmd-intersect}
\cmd{intersection\_check}, Alias \code{intersect}. Berechnet alle Schnitte zwischen den
ausgewählten Objekten und legt sie als neue Objekte an:
\begin{itemize}
\item Kurve mit Kurve und Kurve mit Fläche: Schnittpunkte als Punktobjekte,
\item Fläche mit Fläche: Schnittkurven (und gegebenenfalls einzelne Berührpunkte).
\end{itemize}
Flächen desselben Körpers werden nicht miteinander geschnitten. Ergebnis in der Statusleiste:
\code{Intersections: 1 curve(s), 2 point(s)}. Die erzeugten Punkte und Kurven können Sie mit
dem Fang weiterverwenden, etwa als Schnittkante einer Linse mit einer Blendenebene.

% ============================================================================
\section{Arbeitsabläufe an Beispielen}\label{sec:mod-beispiele}

Die folgenden Beispiele führen Schritt für Schritt durch typische Aufgaben. Sie verwenden
eingetippte Koordinaten, damit Sie die Ergebnisse genau nachvollziehen können; in der Praxis
können Sie jeden Punkt ebenso mit der Maus und dem Fang setzen. \enquote{(Enter)} bedeutet
jeweils: mit \keys{Enter} bestätigen. Die optische Achse ist in allen Beispielen die $z$-Achse.

\subsection{Plankonvexe Linse als Volumenkörper}\label{sec:bsp-linse}
\index{Linse modellieren}

Ziel: eine plankonvexe Linse mit \SI{25}{mm} Durchmesser, \SI{5}{mm} Mittendicke und
\SI{50}{mm} Krümmungsradius, die plane Seite bei $z=0$, der Scheitel bei $z=5$.

Idee: Die Linse ist die Schnittmenge einer Kugel (Radius 50, Mittelpunkt auf der Achse bei
$z = 5 - 50 = -45$) mit einem Zylinder (Radius 12.5, von $z=0$ bis $z=5$).

\begin{schritte}
\item Bewegen Sie die Maus über die Ansicht \textbf{Top}, damit Zylinder in $z$-Richtung
  stehen.
\item Kugel: \code{sph} (Enter), \code{0,0,-45} (Enter), \code{50} (Enter).
\item Zylinder: \code{cyl} (Enter), \code{0,0,0} (Enter), \code{12.5} (Enter), \code{5}
  (Enter).
\item Wählen Sie die Kugel und mit \keys{Shift}+Klick den Zylinder aus (die Reihenfolge ist bei
  der Schnittmenge gleichgültig).
\item \code{bi} (Enter). Kugel und Zylinder werden durch die Linse ersetzt.
\item Kontrolle: Linse auswählen, \code{volume} (Enter). Die Statusleiste zeigt ein Volumen von
  etwa \num{2066.78} (\si{\cubic\milli\metre}). Die Randdicke können Sie mit \gui{Distance}
  zwischen den beiden Randkanten messen (Fang \gui{Corner} bzw.\ \gui{Quad} mit eingeschaltetem
  \gui{Edge}); sie beträgt $5 - (50 - \sqrt{50^2 - 12.5^2}) \approx \num{3.412}$.
\item Mit \gui{Curvature} auf der gewölbten Fläche erhalten Sie die Hauptkrümmungen
  $k_1 = k_2 = \pm 0.02$, also den Radius 50.
\end{schritte}

\paragraph{Bikonvexe Linse.} Die Schnittmenge wird über alle ausgewählten Körper gebildet; für
eine symmetrische Bikonvexlinse mit \SI{6}{mm} Mittendicke genügen daher zwei Kugeln und ein
Zylinder: Kugeln mit Radius 50 um \code{0,0,-47} und \code{0,0,47}, Zylinder um
\code{0,0,-3} mit Radius 12.5 und Höhe 6; alle drei auswählen, \code{bi}.

\subsection{Sphärischer Spiegel als Rotationsfläche}\label{sec:bsp-spiegel}
\index{Spiegel modellieren}

Ziel: ein Hohlspiegel mit Radius \SI{100}{mm} und \SI{50}{mm} Durchmesser, Scheitel im
Ursprung, Krümmungsmittelpunkt bei $(0,0,100)$. Wir zeichnen das halbe Profil als Bogen in der
XZ-Ebene und drehen es um die $z$-Achse.

Die Pfeilhöhe am Rand ($r=25$) ist $100 - \sqrt{100^2 - 25^2} \approx \num{3.1754}$, bei $r =
12.5$ beträgt sie $\approx \num{0.7843}$.

\begin{schritte}
\item Bogen: \code{a} (Enter), dann \code{0,0,0} (Enter) als Anfang, \code{25,0,3.1754}
  (Enter) als Ende und \code{12.5,0,0.7843} (Enter) als Punkt auf dem Bogen. Der Bogen geht
  exakt durch diese drei Punkte; sein Radius ist durch die gerundete Eingabe sehr nahe an 100.
\item Wählen Sie den Bogen aus.
\item \code{rev} (Enter), Achsanfang \code{0,0,0} (Enter), Achsende \code{0,0,10} (Enter),
  dann \keys{Enter} für \SI{360}{\degree}.
\item Prüfen Sie mit \gui{Curvature} auf der Fläche den Radius.
\item Optional eine Rückseite: Fläche auswählen, \code{offsetsrf} (Enter), \code{5} oder
  \code{-5} (Enter) -- das Vorzeichen wählt die Seite. Prüfen Sie in der Ansicht Front, ob die
  Parallelfläche hinter dem Spiegel liegt; sonst \keys{Strg+Z} und das andere Vorzeichen.
\end{schritte}

Für asphärische oder freigeformte Profile zeichnen Sie statt des Bogens eine
\gui{Interpolated spline} durch berechnete Stützpunkte und verfahren ebenso.

\subsection{Prisma aus einem Profil: Extrudieren und Deckeln}\label{sec:bsp-prisma}

Ziel: ein sechseckiger Stab mit Umkreisradius 10 und Länge 20 als geschlossener Körper.

\begin{schritte}
\item Maus über \textbf{Top}. Polygon: \code{pol} (Enter), \code{0,0} (Enter), \code{10}
  (Enter). Es entsteht ein Sechseck in der XY-Ebene.
\item Sechseck auswählen, \code{ext} (Enter), \code{20} (Enter). Es entsteht der offene
  Mantel.
\item Mantel auswählen, \code{cap} (Enter). Die Statusleiste meldet
  \code{Capped: closed solid, volume 5196.1524} -- das ist
  $\frac{3\sqrt{3}}{2} \cdot 10^2 \cdot 20$.
\end{schritte}

Auf dieselbe Weise entstehen aus einem Rechteck eine Platte, aus einem Kreis ein Zylinder und
aus einer geschlossenen Polylinie beliebige prismatische Körper, etwa Halterungen oder
Blendenträger.

\subsection{Platte mit Bohrung: Boolesche Differenz}\label{sec:bsp-platte}

Ziel: eine Blende -- eine Platte $40 \times 40 \times 5$ mit einer zentralen Bohrung von
\SI{10}{mm} Durchmesser.

\begin{schritte}
\item Maus über \textbf{Top}. Quader: \code{b} (Enter), \code{-20,-20} (Enter),
  \code{20,20} (Enter), \code{5} (Enter).
\item Zylinder, der oben und unten übersteht: \code{cyl} (Enter), \code{0,0,-1} (Enter),
  \code{5} (Enter), \code{7} (Enter).
\item Klicken Sie \emph{zuerst} auf die Platte, dann mit \keys{Shift}+Klick auf den Zylinder.
\item \code{bd} (Enter). Die Platte hat jetzt eine Bohrung.
\item Kontrolle mit \code{volume}: $40 \cdot 40 \cdot 5 - \pi \cdot 5^2 \cdot 5 \approx
  \num{7607.30}$.
\end{schritte}

Für mehrere Bohrungen auf einem Lochkreis setzen Sie einen Zylinder außerhalb der Mitte, wählen
ihn aus und vervielfältigen ihn mit \gui{Polar array} um den Ursprung, bevor Sie Platte und
alle Zylinder für die Differenz auswählen (Platte zuerst).

\subsection{Flächen aus Kurven: Loft und Sweep}\label{sec:bsp-loft}

\paragraph{Konischer Übergang mit Loft.}
\begin{schritte}
\item Maus über \textbf{Top}. Kreis \code{c} (Enter), \code{0,0,0} (Enter), \code{10} (Enter).
\item Zweiter Kreis in der Höhe 20: \code{c} (Enter), \code{0,0,20} (Enter), \code{5} (Enter).
\item Ersten Kreis anklicken, zweiten mit \keys{Shift}+Klick dazunehmen.
\item \code{loft} (Enter). Es entsteht die Mantelfläche eines Kegelstumpfs; mit \code{cap}
  schließen Sie sie zu einem Körper.
\end{schritte}
Mit drei oder mehr Querschnitten (in Auswahlreihenfolge) entstehen gewölbte Übergänge, etwa
für Lichtleiter oder Trichter.

\paragraph{Rohr mit Sweep 1 rail.}
\begin{schritte}
\item Zeichnen Sie die Leitkurve, z.\,B.\ eine \gui{Interpolated spline} oder einen Bogen.
\item Zeichnen Sie das Profil, einen kleinen Kreis, an den Anfang der Leitkurve -- in der
  Ansicht, deren Konstruktionsebene senkrecht zur Leitkurve an ihrem Anfang steht; den
  Mittelpunkt fangen Sie mit \gui{End}.
\item Profil anklicken, Leitkurve mit \keys{Shift}+Klick dazunehmen, \code{sweep1} (Enter).
\end{schritte}

\subsection{Linsen-Array mit Grid array}\label{sec:bsp-array}

\begin{schritte}
\item Erzeugen Sie eine Linse wie in Abschnitt~\ref{sec:bsp-linse} und wählen Sie sie aus.
\item Maus über \textbf{Top}. \code{array\_grid} (Enter).
\item \code{5} (Enter), \code{4} (Enter) -- die Aufforderung zeigt jetzt
  \code{Grid array (5 $\times$ 4): base point \ldots}.
\item Basispunkt \code{0,0} (Enter), erster Schritt \code{@30,0} (Enter), zweiter Schritt
  \code{@0,30} (Enter).
\end{schritte}
Es entstehen 19 Kopien; zusammen mit dem Original sind es $5 \times 4 = 20$ Linsen im Abstand
von \SI{30}{mm}.

\subsection{Kurven stutzen und verbinden}\label{sec:bsp-trim}

Ein häufiger Ablauf beim Zeichnen von Profilen: Linien grob über ihre Schnittpunkte hinaus
zeichnen, dann die Überstände entfernen und alles zu einer geschlossenen Kurve verbinden.

\begin{schritte}
\item Zeichnen Sie zwei sich kreuzende Linien.
\item Wählen Sie die erste Linie (die gekürzt werden soll), dann mit \keys{Shift}+Klick die
  zweite (das Schneidobjekt).
\item \code{tr} (Enter) und klicken Sie auf den überstehenden Teil der ersten Linie. Er
  verschwindet; der Rest bleibt als neue Linie.
\item Wiederholen Sie das umgekehrt für die zweite Linie.
\item Wählen Sie alle Teilkurven des Profils aus und verbinden Sie sie mit \code{j} (Enter) zu
  einer Kurve. Ecken runden Sie anschließend mit \code{fillet} ab (Kurve auswählen, Befehl
  starten, Radius tippen).
\end{schritte}

\begin{tipp}
Nach jedem Schritt können Sie mit \keys{Strg+Z} zurückgehen. Das Menü \menu{Edit} zeigt, welcher
Befehl zuletzt ausgeführt wurde. So probieren Sie gefahrlos Varianten aus -- etwa das
Vorzeichen einer Offset-Distanz oder die Reihenfolge der Auswahl.
\end{tipp}

\section{Häufige Fehler}\label{sec:modellieren-fehler}

\begin{referenz}{p{0.38\linewidth}X}
\toprule
Problem & Ursache und Abhilfe \\
\midrule\endhead
Der Befehl startet nicht, die Statusleiste meldet \enquote{select \ldots} &
  Der Befehl arbeitet auf der Auswahl. Wählen Sie die Objekte \emph{vor} dem Start aus
  (Abschnitt~\ref{sec:auswahlbefehle}). \\
Ein Kreis, Rechteck oder Zylinder liegt in der falschen Ebene &
  Die Maus war zuletzt über einer anderen Ansicht. Bewegen Sie die Maus vor dem Tippen über die
  gewünschte Ansicht oder setzen Sie den ersten Punkt dort mit der Maus. \\
Eingetippte Koordinaten landen woanders als erwartet &
  Mit zwei Werten gelten Ebenenkoordinaten der Ansicht (in Front ist der zweite Wert $z$).
  Verwenden Sie drei Werte für Weltkoordinaten. Dezimaltrennzeichen ist der Punkt. \\
Der Punkt springt an eine unerwartete Stelle &
  Der Objektfang hat gegriffen. Prüfen Sie das Fangsymbol; \keys{Tab} wechselt den Kandidaten,
  \keys{Alt} unterdrückt den Fang, \keys{F3} schaltet ihn ab. \\
Ortho wirkt nicht &
  Gefangene und eingetippte Punkte werden von Ortho nicht verändert (Abschnitt~\ref{sec:ortho}). \\
Boolesche Operation meldet \enquote{Select at least two solids} oder \enquote{Booleans need
closed solids} &
  Ein Objekt ist kein geschlossener Körper (z.\,B.\ eine ungedeckelte Extrusion oder eine offene
  Hülle). Schließen Sie es mit Cap oder Join; prüfen Sie mit Volume. \\
Boolean difference zieht das Falsche ab &
  Die Auswahlreihenfolge war vertauscht. \keys{Strg+Z}, dann zuerst den Grundkörper wählen. \\
Offset surface oder Fillet surfaces nehmen eine Fläche nicht an &
  Es ist eine Fläche eines Körpers oder eine beschnittene Fläche. Diese Befehle arbeiten nur mit
  unbeschnittenen Einzelflächen. \\
Loft verdreht die Fläche &
  Die Querschnitte laufen in unterschiedlicher Richtung oder wurden in falscher Reihenfolge
  gewählt. \\
\code{h} startet nicht die Helix &
  In einer Ansicht blendet \keys{H} die Auswahl aus, solange etwas ausgewählt ist. Heben Sie
  die Auswahl mit \keys{Esc} auf, oder tippen Sie in der Befehlszeile; alternativ starten Sie
  die Helix über die Werkzeugspalte. \\
Eine eingetippte Extrude-Distanz geht in die falsche Richtung &
  Sie wirkt entlang der Normalen der Ansicht, über der die Maus zuletzt war (Front: $-y$, Left:
  $-x$, Top: $+z$). Bewegen Sie die Maus vor dem Tippen über die gewünschte Ansicht, ändern Sie
  das Vorzeichen oder geben Sie die Richtung mit zwei Punkten an. \\
Ein Rechtsklick startet unerwartet einen Befehl &
  Ohne laufenden Befehl wiederholt ein Rechtsklick (ohne Ziehen) in einer Ansicht den letzten
  Befehl. Zum Drehen oder Verschieben ziehen Sie mit der rechten Taste. \keys{Esc} bricht den
  wiederholten Befehl ab; hat ein Auswahlbefehl (etwa eine Boolesche Operation) sofort ein
  Ergebnis erzeugt, nehmen Sie es mit \keys{Strg+Z} zurück. \\
\bottomrule
\end{referenz}

\fussy
